\section{Statistik und Statistical Learning 2}

\subsection{Skalenniveaus und Charakterisierung von Merkmalen}

\textbf{Datentypen:} \textbf{Strukturierte Daten} (vordefinierte Form, Tabellen); \textbf{Unstrukturierte Daten} (Texte, Bilder, Audiodateien, Sensordaten).

\definition{\textbf{Datenmatrix:} $n \times p$-Matrix mit Merkmalsträger $i=1,\ldots,n$ (Zeilen, Objekte) und Merkmale $j=1,\ldots,p$ (Spalten, Variablen). Merkmalsausprägungen sind die Werte in dieser Matrix.}

\definition{\textbf{Klassifikation von Variablen:}}
\begin{center}
    \small
    \begin{tabular}{|l|l|}
        \hline
        \textbf{Qualitativ (kategorial)} & Endliche Kategorien (z.\,B.\ Geschlecht, Augenfarbe)                         \\
        \hline
        \textbf{Quantitativ (numerisch)} & Zahlenwerte (z.\,B.\ Alter, Größe)                                          \\
        \hline
        \textbf{Diskret}                 & Endlich/abzählbar viele Ausprägungen (z.\,B.\ Anzahl)                        \\
        \hline
        \textbf{Stetig}                  & Unendlich viele Werte im Intervall (z.\,B.\ Gewicht)                         \\
        \hline
        \textbf{Manifest}                & Direkt beobachtbar (z.\,B.\ Körpergröße)                                      \\
        \hline
        \textbf{Latent}                  & Nicht direkt beobachtbar, durch Indikatoren gemessen (z.\,B.\ Zufriedenheit) \\
        \hline
        \textbf{Unabhängig}              & Einflussgröße, Prädiktor                                                     \\
        \hline
        \textbf{Abhängig}                & Zielgröße, Kriterium                                                          \\
        \hline
    \end{tabular}
\end{center}

\definition{\textbf{Skalenniveaus (Messniveaus):} Hierarchische Klassifikation für zulässige Operationen:}
\begin{center}
    \small
    \begin{tabular}{|l|l|l|}
        \hline
        \textbf{Skala} & \textbf{Operationen}      & \textbf{Lagemaß}          \\
        \hline
        Nominal        & $=$, $\neq$               & Modus                     \\
        \hline
        Ordinal        & $=$, $<$, $>$             & Modus, Median             \\
        \hline
        Intervall      & $=$, $<$, $+$, $-$        & Modus, Median, Mittelwert \\
        \hline
        Verhältnis     & alle                      & Modus, Median, Mittelwert \\
        \hline
        Absolut        & alle (natürliche Einheit) & Modus, Median, Mittelwert \\
        \hline
    \end{tabular}
\end{center}

\newpage

\subsection{Grafische Darstellung und Visualisierung}

\definition{\textbf{Univariate Darstellung:}}
\begin{center}
    \small
    \begin{tabular}{|l|l|}
        \hline
        \textbf{Diagrammtyp} & \textbf{Einsatz}                              \\
        \hline
        Balkendiagramm       & Kategoriale/metrische Häufigkeiten           \\
        \hline
        Histogramm           & Kontinuierliche Verteilungen (Klassen)       \\
        \hline
        Boxplot              & Lage, Streuung, Ausreißer (Median, Quartile) \\
        \hline
        Violin-Plot          & Verteilungsform + Boxplot-Info               \\
        \hline
    \end{tabular}
\end{center}

\definition{\textbf{Bivariate Darstellung:}}
\begin{center}
    \small
    \begin{tabular}{|l|l|}
        \hline
        \textbf{Diagrammtyp}        & \textbf{Einsatz}                       \\
        \hline
        Streudiagramm (Scatterplot) & 2 metrische Variablen, Korrelation     \\
        \hline
        Boxplot nach Gruppen        & Metrisch vs. Kategorial, Vergleich     \\
        \hline
        Gruppiertes Balkendiagramm  & 2 kategoriale oder gemischte Variablen \\
        \hline
        Liniendiagramm              & Zeitreihen, Trends                     \\
        \hline
        Heatmap                     & Korrelationsmatrizen, 2D-Intensitäten  \\
        \hline
    \end{tabular}
\end{center}

\newpage

\subsection{Ähnlichkeitsmaße und Distanzmaße}

\definition{\textbf{Grundkonzept:} Distanz $d$ misst Unähnlichkeit; Ähnlichkeit $s$ und Distanz sind invers: $s = 1 - d$ (normalisiert).}

\definition{\textbf{Gängige Distanzmaße:}}
\begin{itemize}[leftmargin=1.5em]
    \item \textbf{Euklidisch:} $d = \sqrt{\sum_{j=1}^{p}(x_{B,j}-x_{A,j})^2}$ - Luftlinienabstand; empfindlich für Ausreißer.
    \item \textbf{Manhattan:} $d = \sum_{j=1}^{p}|x_{B,j}-x_{A,j}|$ - robuster gegenüber Ausreißern.
    \item \textbf{Minkowski:} $d_k = \left(\sum_{j=1}^{p}|x_{B,j}-x_{A,j}|^k\right)^{1/k}$ - Verallgemeinerung ($k$-Parameter).
    \item \textbf{Mahalanobis:} $d = \sqrt{(\mathbf{x}_B-\mathbf{x}_A)^T\hat{\boldsymbol{\Sigma}}^{-1}(\mathbf{x}_B-\mathbf{x}_A)}$ - berücksichtigt Kovarianzstruktur.
\end{itemize}

\definition{\textbf{Ähnlichkeitsmaße:}}
\begin{itemize}[leftmargin=1.5em]
    \item \textbf{Kosinus-Ähnlichkeit:} $\cos(\theta) = \frac{\mathbf{x}_A \cdot \mathbf{x}_B}{\|\mathbf{x}_A\|_2\|\mathbf{x}_B\|_2}$ - Winkel zwischen Vektoren (Text, Bilder).
    \item \textbf{Tanimoto (Jaccard):} $S = \frac{|A \cap B|}{|A \cup B|}$ - für Binärvektoren, $S \in [0,1]$.
    \item \textbf{Dice:} $S = \frac{2|A \cap B|}{|A|+|B|}$ - ähnlich Tanimoto, stärker gewichtet.
    \item \textbf{Pearson-/Spearman-Distanz:} $d = 1-r$ - korrelationsbasiert, robust (Spearman).
\end{itemize}

\textbf{Wahl des Maßes:} Euklidisch (normalisierte numerische Daten) • Mahalanobis (korrelierte Variablen) • Kosinus (hochdimensional: Text, Bilder) • Tanimoto/Dice (binär) • Spearman (Ausreißer, Ränge) • Manhattan (Robustheit)

\newpage

\subsection{Wahrscheinlichkeiten und Wahrscheinlichkeitsräume}

\definition{\textbf{Zufallsexperiment:} Prozess mit unbekanntem Ausgang, aber bekannten möglichen Ergebnissen (Elementarereignisse $\omega$).}

\textbf{Wahrscheinlichkeitsinterpretationen:}
\textbf{Frequentistisch} - relative Häufigkeit bei sehr häufiger Wiederholung (objektiv). • 
\textbf{Bayessch} - subjektiver Grad der Gewissheit basierend auf verfügbarem Wissen.

\definition{\textbf{Wahrscheinlichkeitsraum $(\Omega, \mathcal{A}, P)$:} $\Omega$ = Stichprobenraum (alle möglichen Ausgänge); $\mathcal{A}$ = $\sigma$-Algebra (Menge aller Ereignisse mit zugeordneten WKen); $P$ = Wahrscheinlichkeitsfunktion $P: \mathcal{A} \to [0,1]$.}

\textbf{Kolmogorov-Axiome:}
(1) $P(\Omega) = 1$; 
(2) $P(E) \geq 0$ für alle $E \in \mathcal{A}$; 
(3) Additivität disjunkter Ereignisse: $P(E_1 \cup E_2) = P(E_1) + P(E_2)$ für $E_1 \cap E_2 = \emptyset$.

\definition{\textbf{Laplace'scher Wahrscheinlichkeitsraum:} Endlicher Raum mit gleichen Elementarereignis-WKen. Für Ereignis $E$ mit $g$ günstigen Fällen und $m$ möglichen Fällen: $P(E) = \frac{g}{m}$.}

\textbf{Abgeleitete Regeln:}
Komplementär: $P(E^c) = 1-P(E)$ • 
Addition: $P(E_1 \cup E_2) = P(E_1) + P(E_2) - P(E_1 \cap E_2)$ • 
Monotonie: $E_1 \subseteq E_2 \Rightarrow P(E_1) \leq P(E_2)$

\newpage





\subsection{Kombinatorik}

\definition{\textbf{Kombinatorik:} Zählen von Anordnungen und Auswahlen endlicher Mengen. Fundamental für Wahrscheinlichkeitsberechnungen: $P(E) = \frac{\text{günstige Fälle}}{\text{mögliche Fälle}}$.}

\definition{\textbf{Kombinatorische Formeln:}}
\begin{center}
    \small
    \begin{tabular}{|l|l|l|}
        \hline
        \textbf{Typ}                          & \textbf{Ohne Wiederholung}             & \textbf{Mit Wiederholung}     \\
        \hline
        Permutation ($n$ Elemente anordnen)   & $n!$                                   & $\dfrac{n!}{n_1!\cdots n_k!}$ \\[8pt]
        \hline
        Variation ($k$ aus $n$, geordnet)     & $\dfrac{n!}{(n-k)!}$                   & $n^k$                         \\[8pt]
        \hline
        Kombination ($k$ aus $n$, ungeordnet) & $\dbinom{n}{k} = \dfrac{n!}{k!(n-k)!}$ & $\dbinom{n+k-1}{k}$           \\[8pt]
        \hline
    \end{tabular}
\end{center}

\textbf{Formelwahl:} Alle $n$ Elemente? $\to$ Permutation. Nur $k < n$ ausgewählt? Reihenfolge wichtig? $\to$ Variation. Nicht? $\to$ Kombination. Wiederholungen erlaubt? $\to$ entsprechende Spalte.

\subsection{Diskrete und stetige Zufallsvariablen}
\newpage

\subsection{Statistische Tests}

\definition{\textbf{Hypothesentest:} Verfahren, um anhand von Stichprobendaten zu entscheiden, ob eine Vermutung über einen Parameter oder eine Verteilung haltbar ist.}

\definition{\textbf{Hypothesen:} $H_0$ (Nullhypothese) = konservative Annahme, meist ``kein Effekt''. $H_1$ (Alternativhypothese) = Gegenhypothese. Unterscheidung: \emph{einfach} (einzelner Wert) vs. \emph{zusammengesetzt} (Wertebereich); \emph{einseitig} ($H_1: \mu > \mu_0$) vs. \emph{zweiseitig} ($H_1: \mu \neq \mu_0$).}

\definition{\textbf{Signifikanzniveau $\alpha$:} Maximal akzeptierte Wahrscheinlichkeit für Fehler 1. Art. Typisch: $\alpha \in \{0.01, 0.05, 0.10\}$.}

\textbf{Ablauf des Tests:}
(1) $H_0, H_1, \alpha$ formulieren • (2) Teststatistik $T$ berechnen • (3) Kritischer Bereich bestimmen • (4) Entscheidung treffen.

\definition{\textbf{p-Wert:} Wahrscheinlichkeit, unter $H_0$ einen Wert $\geq t_{\text{obs}}$ zu beobachten (in Richtung $H_1$). Entscheidung: $p < \alpha \Rightarrow H_0$ ablehnen; $p \geq \alpha \Rightarrow H_0$ nicht ablehnen.}

\definition{\textbf{Testfehler:}}
\begin{center}
    \small
    \begin{tabular}{|l|c|c|}
        \hline
                             & $H_0$ wahr                     & $H_0$ falsch                 \\
        \hline
        $H_0$ nicht ablehnen & Korrekt (Konfidenz $1-\alpha$) & \textbf{Fehler II} ($\beta$) \\
        \hline
        $H_0$ ablehnen       & \textbf{Fehler I} ($\alpha$)   & Korrekt (Power $1-\beta$)    \\
        \hline
    \end{tabular}
\end{center}

Fehler I (Fehlalarm): $H_0$ ablehnen, obwohl wahr, $P = \alpha$ • 
Fehler II (falsch negativ): $H_0$ nicht ablehnen, obwohl falsch, $P = \beta$ • 
Power = $1 - \beta$ (Wahrscheinlichkeit, $H_0$ richtig abzulehnen).

Tradeoff: $\alpha$ und $\beta$ negativ korreliert; Reduktion beider erfordert größere Stichprobe.
\newpage

\subsection{Population, Stichproben und Konfidenzintervalle}

\definition{\textbf{Grundkonzept:} Population = alle Objekte (unbekannte Verteilung); Stichprobe = Teilmenge (Zufallsstichprobe). Inferenzprinzip: Empirische Stichprobenverteilung $\to$ wahre Populationsverteilung.}

\definition{\textbf{Schätzer:}}
\begin{center}
    \small
    \begin{tabular}{|l|l|}
        \hline
        \textbf{Parameter} & \textbf{Schätzer}                           \\
        \hline
        Mittelwert $\mu$   & $\bar{X} = \frac{1}{n}\sum X_i$             \\
        \hline
        Varianz $\sigma^2$ & $S^2 = \frac{1}{n-1}\sum (X_i - \bar{X})^2$ \\
        \hline
        Anteil $p$         & $\hat{p} = \frac{\text{Erfolge}}{n}$        \\
        \hline
    \end{tabular}
\end{center}

\definition{\textbf{Konfidenzintervall:} Zufälliges Intervall $[\theta_l, \theta_u]$ mit $P(\theta_l \leq \theta \leq \theta_u) = 1-\alpha$. Das Intervall ist zufällig (hängt von Stichprobe ab), der Parameter $\theta$ ist fest. Interpretation: Bei Wiederholung liegt der wahre Parameter in $(1-\alpha) \times 100\%$ aller KIs. Konfidenzniveau typisch $1-\alpha \in \{0.90, 0.95, 0.99\}$. \textbf{KI-Typen:} Zweiseitig $[\theta_l, \theta_u]$ mit $P(\theta < \theta_l) = P(\theta > \theta_u) = \frac{\alpha}{2}$ • Einseitig unten $[\theta_l, \infty)$ • Einseitig oben $(-\infty, \theta_u]$.}

\definition{\textbf{KI für $\mu$ (bekannte $\sigma_0^2$):} Voraussetzung: Population normal oder $n \geq 30$ (CLT). Teststatistik: $Z = \frac{\bar{X} - \mu}{\sigma_0/\sqrt{n}} \sim N(0,1)$. KI: $\bar{X} \pm z_{1-\alpha/2} \frac{\sigma_0}{\sqrt{n}}$ (for $\alpha=0.05$: $z \approx 1.96$).}

\definition{\textbf{KI für $\mu$ (unbekannte $\sigma^2$):} Teststatistik: $T = \frac{\bar{X} - \mu}{S/\sqrt{n}} \sim t_{n-1}$ (mit $S^2 = \frac{1}{n-1}\sum(X_i-\bar{X})^2$). KI: $\bar{X} \pm t_{n-1, 1-\alpha/2} \frac{S}{\sqrt{n}}$. Die $t$-Verteilung hat schwerere Schwänze als Normal (besonders kleines $n$) $\to$ breitere KIs.}

\definition{\textbf{KI für $\sigma^2$ (normalverteilte Pop.):} Teststatistik: $Q = \frac{(n-1)S^2}{\sigma^2} \sim \chi^2_{n-1}$. KI: $\left[\frac{(n-1)S^2}{\chi^2_{n-1, 1-\alpha/2}}, \frac{(n-1)S^2}{\chi^2_{n-1, \alpha/2}}\right]$.}

\subsection{Zentraler Grenzwertsatz und Gesetz der großen Zahlen}

\textbf{Zwei fundamentale Theoreme der Wahrscheinlichkeitstheorie:}

\definition{\textbf{Unabhängig und identisch verteilt (u.i.v., iid):} Wenn ein Zufallsexperiment — dargestellt durch eine Zufallsvariable $X$ mit Erwartungswert $\mu$, Varianz $\sigma^2$ und Verteilungsfunktion $F(x)$ — unabhängig $n$-mal wiederholt wird, so bezeichnen wir das Ergebnis beim $i$-ten Versuch mit $X_i$ für $i = 1, \ldots, n$. Die Zufallsvariablen $X_1, \ldots, X_n$ sind unabhängig voneinander und identisch verteilt (haben dieselbe Verteilung wie $X$), folglich haben alle denselben Erwartungswert $\mu$ und dieselbe Varianz $\sigma^2$. Man kann $X_1, \ldots, X_n$ als unabhängige ``Kopien'' von $X$ auffassen. Die konkreten Zahlenwerte $x_1, \ldots, x_n$ heißen Realisierungen.}

\definition{\textbf{Markov-Ungleichung:} Für nicht-negative $X$ und $a > 0$: $P(X \geq a) \leq \frac{E[X]}{a}$. Modellunabhängige obere Schranke.}

\definition{\textbf{Tschebyscheff-Ungleichung:} Für beliebiges $X$ mit $E[X]=\mu$, $\text{Var}(X)=\sigma^2$ und $c > 0$: $P(|X-\mu| \geq c) \leq \frac{\sigma^2}{c^2}$. Vereinfacht: $P(|X-\mu| \geq k\sigma) \leq \frac{1}{k^2}$, z.B. $k=2$: mind. 75\% im Intervall $[\mu-2\sigma, \mu+2\sigma]$; $k=3$: mind. 89\%.}

\definition{\textbf{Gesetz der großen Zahlen (LLN):} Für u.i.v. $X_1, \ldots, X_n$ mit $E[X]=\mu$, $\text{Var}(X)=\sigma^2$ gilt: $P(|\bar{X}_n - \mu| < c) \to 1$ für $n \to \infty$ (für beliebig kleines $c>0$). D.h. $\bar{X}_n \xrightarrow{P} \mu$. Intuition: Große Stichproben liefern vertrauenswürdige Schätzungen. Bernoulli-Spezialfall: Relative Häufigkeit konvergiert zur Wahrscheinlichkeit.}

\definition{\textbf{Zentraler Grenzwertsatz (CLT):} Für u.i.v. $X_1, \ldots, X_n$ mit $E[X]=\mu$, $\text{Var}(X)=\sigma^2$ ist die Summe asymptotisch normalverteilt: $S_n \xrightarrow{d} N(n\mu, n\sigma^2)$. Standardisierte Form: $Z_n = \frac{\bar{X}_n - \mu}{\sigma/\sqrt{n}} \xrightarrow{d} N(0,1)$. Das Besondere: Dies gilt \emph{unabhängig} von der Verteilung der $X_i$ (gleichmäßig, exponentiell, etc.). Faustregel: wirkt ab $n \approx 30$. Praktische Relevanz: Erklärt, warum viele Phänomene normalverteilt sind (Summen vieler Einflussfaktoren); rechtfertigt KI-Konstruktion und $z$-Tests.}

\newpage

\subsection{Mehrdimensionale Zufallsvariablen}

\definition{\textbf{Mehrdimensionale (Bivariate) Zufallsvariablen:} Ein zufälliger Vektor oder eine gemeinsame Zufallsvariable mit $n$ Komponenten:
\formel{\mathbf{X} = \begin{pmatrix} X_1 \\ X_2 \\ \vdots \\ X_n \end{pmatrix} \quad \text{oder im bivariaten Fall:} \quad (X, Y) = \begin{pmatrix} X \\ Y \end{pmatrix}}
wobei jede Komponente selbst eine diskrete oder stetige Zufallsvariable ist. Die Wertemenge ist eine Teilmenge von $\mathbb{R}^n$. \textbf{Beispiel:} Bei einer zufällig ausgewählten Person interessieren wir uns für $(X, Y)$ mit $X$ = Körpergröße, $Y$ = Körpergewicht.}

\definition{\textbf{Gemeinsame Verteilung (Joint Distribution):}
\textbf{Im diskreten Fall:} Die gemeinsame Wahrscheinlichkeitsfunktion (joint PMF) ist:
\formel{f_{X,Y}(x, y) = P(X = x, Y = y)}
mit der Normierungsbedingung: $\sum_x \sum_y f_{X,Y}(x, y) = 1$.

\textbf{Im stetigen Fall:} Die gemeinsame Dichtefunktion (joint PDF) ist:
\formel{f_{X,Y}(x, y) \geq 0 \quad \text{mit} \quad \int_{-\infty}^{\infty} \int_{-\infty}^{\infty} f_{X,Y}(x, y) \, dx \, dy = 1}

Die gemeinsame Verteilungsfunktion (CDF) für beide Fälle:
\formel{F_{X,Y}(x, y) = P(X \leq x, Y \leq y)}
Die gemeinsame Verteilung enthält vollständige Information über das Verhalten beider Variablen und deren Abhängigkeitsstruktur.}

\definition{\textbf{Randverteilungen (Marginal Distributions):} Die Verteilung einer einzelnen Komponente, wenn wir die Information über die anderen ``heraussummieren'':
\textbf{Im diskreten Fall:}
\formel{f_X(x) = \sum_y f_{X,Y}(x, y) \quad \text{und} \quad f_Y(y) = \sum_x f_{X,Y}(x, y)}
\textbf{Im stetigen Fall:}
\formel{f_X(x) = \int_{-\infty}^{\infty} f_{X,Y}(x, y) \, dy \quad \text{und} \quad f_Y(y) = \int_{-\infty}^{\infty} f_{X,Y}(x, y) \, dx}
Die Randverteilungen sind wieder Verteilungen eindimensionaler Zufallsvariablen. Sie enthalten weniger Information als die gemeinsame Verteilung. \textbf{Beispiel (Kontingenztafel):} Bei diskreten Variablen sind die Randverteilungen die Zeilen- und Spaltensummen der Kontingenztafel.}

\definition{\textbf{Bedingte Verteilungen (Conditional Distributions):} Die Verteilung einer Variablen unter der Bedingung, dass eine andere einen bestimmten Wert hat.
\textbf{Im diskreten Fall:}
\formel{P(X = x \mid Y = y) = \frac{P(X = x, Y = y)}{P(Y = y)} = \frac{f_{X,Y}(x, y)}{f_Y(y)}}
wenn $P(Y = y) > 0$.
\textbf{Im stetigen Fall:}
\formel{f_{X \mid Y}(x \mid y) = \frac{f_{X,Y}(x, y)}{f_Y(y)}}
wenn $f_Y(y) > 0$. Die bedingte Dichtefunktion erfüllt: $\int_{-\infty}^{\infty} f_{X \mid Y}(x \mid y) \, dx = 1$ für jedes feste $y$.}

\definition{\textbf{Unabhängigkeit von Zufallsvariablen:} Zwei Zufallsvariablen $X$ und $Y$ sind unabhängig, wenn ihre gemeinsame Verteilung das Produkt der Randverteilungen ist:
\textbf{Im diskreten Fall:}
\formel{f_{X,Y}(x, y) = f_X(x) \cdot f_Y(y) \quad \text{für alle } x, y}
\textbf{Im stetigen Fall:}
\formel{f_{X,Y}(x, y) = f_X(x) \cdot f_Y(y) \quad \text{für alle } x, y}
\textbf{Äquivalente Bedingung (für bedingte Wahrscheinlichkeiten):} Die bedingte Verteilung gleicht der Randverteilung:
\formel{P(X = x \mid Y = y) = P(X = x) \quad \text{und} \quad P(Y = y \mid X = x) = P(Y = y)}
\textbf{Intuition:} Die Information, dass $Y$ einen bestimmten Wert hat, ändert unsere Wahrscheinlichkeitsaussagen über $X$ nicht. \textbf{Wichtig:} Unkorreliertheit $\Rightarrow$ NICHT notwendig unabhängig, aber Unabhängigkeit $\Rightarrow$ Unkorreliertheit.}

\textbf{Zusammenfassung (Schema):}
\begin{table}[h]
    \centering
    \small
    \begin{tabular}{|l|l|l|}
        \hline
        \textbf{Konzept} & \textbf{Diskret}                                & \textbf{Stetig}                                         \\
        \hline
        Gemeinsam        & $f_{X,Y}(x, y) = P(X=x, Y=y)$                   & $f_{X,Y}(x, y)$ Dichtefunktion                          \\
        \hline
        Rand ($X$)       & $f_X(x) = \sum_y f_{X,Y}(x, y)$                 & $f_X(x) = \int f_{X,Y}(x, y) dy$                        \\
        \hline
        Bedingt          & $P(X=x \mid Y=y) = \frac{f_{X,Y}(x,y)}{f_Y(y)}$ & $f_{X \mid Y}(x \mid y) = \frac{f_{X,Y}(x,y)}{f_Y(y)}$   \\
        \hline
        Unabhängig       & $f_{X,Y}(x,y) = f_X(x) \cdot f_Y(y)$            & $f_{X,Y}(x,y) = f_X(x) \cdot f_Y(y)$                     \\
        \hline
    \end{tabular}
\end{table}

\newpage

\subsection{Kovarianz und Korrelation}

\definition{\textbf{Kovarianz (Covariance):} Ein Maß für den linearen Zusammenhang zwischen zwei Zufallsvariablen $X$ und $Y$. Sie misst, wie zwei Variablen gemeinsam von ihren Mittelwerten abweichen.}

\textbf{Definition:}
\formel{\text{Cov}(X, Y) = \sigma_{XY} = \mathbb{E}[(X - \mu_X)(Y - \mu_Y)] = \mathbb{E}[XY] - \mathbb{E}[X] \mathbb{E}[Y]}

\textbf{Im Stichprobenfall (empirische Kovarianz):}
\formel{s_{XY} = \frac{1}{n-1} \sum_{i=1}^n (x_i - \bar{x})(y_i - \bar{y})}

\textbf{Interpretation:}
\begin{itemize}[leftmargin=1.5em]
    \item $\text{Cov}(X, Y) > 0$: Positive Assoziation (wenn $X$ größer als $\mu_X$, dann tendenziell $Y > \mu_Y$).
    \item $\text{Cov}(X, Y) < 0$: Negative Assoziation (wenn $X$ größer als $\mu_X$, dann tendenziell $Y < \mu_Y$).
    \item $\text{Cov}(X, Y) = 0$: Keine lineare Assoziation (unkorreliert).
\end{itemize}

\textbf{Probleme:} Die Kovarianz hängt von der Skalierung der Variablen ab — zwei äquivalente Variablen mit unterschiedlichen Einheiten haben unterschiedliche Kovarianzen. Dies erschwert Vergleiche.

\definition{\textbf{Korrelationskoeffizient (Pearson's Correlation Coefficient):} Das standardisierte Maß für den linearen Zusammenhang zwischen $X$ und $Y$:}
\formel{\rho_{X,Y} = \frac{\text{Cov}(X, Y)}{\sigma_X \cdot \sigma_Y} \in [-1, 1]}

\textbf{Im Stichprobenfall (empirischer Korrelationskoeffizient):}
\formel{r = \frac{\sum_{i=1}^n (x_i - \bar{x})(y_i - \bar{y})}{\sqrt{\sum_{i=1}^n (x_i - \bar{x})^2} \cdot \sqrt{\sum_{i=1}^n (y_i - \bar{y})^2}} = \frac{s_{XY}}{s_X \cdot s_Y}}

\textbf{Interpretation:}
\begin{itemize}[leftmargin=1.5em]
    \item $\rho = 1$: Perfekte positive lineare Beziehung ($Y = a + bX$ mit $b > 0$).
    \item $\rho = -1$: Perfekte negative lineare Beziehung ($Y = a + bX$ mit $b < 0$).
    \item $\rho = 0$: Keine lineare Beziehung (kann aber nichtlineare Beziehung existieren!).
    \item $|\rho|$ nähe 0: Schwache lineare Beziehung; $|\rho|$ nahe 1: Starke lineare Beziehung.
\end{itemize}

\textbf{Faustregel für Stärke:} $|\rho| > 0{,}7$ = Stark; $0{,}4 < |\rho| \leq 0{,}7$ = Moderat; $|\rho| \leq 0{,}4$ = Schwach (variiert je nach Kontext).

\definition{\textbf{Eigenschaften der Korrelation:}}
\begin{itemize}[leftmargin=1.5em]
    \item \textbf{Dimensionslos:} Im Gegensatz zur Kovarianz ist die Korrelation unabhängig von der Skalierung.
    \item \textbf{Symmetrisch:} $\rho_{X,Y} = \rho_{Y,X}$.
    \item \textbf{Invariant unter linearen Transformationen:} Wenn $X' = aX + b$ und $Y' = cY + d$ mit $a, c > 0$, dann $\text{Corr}(X', Y') = \text{Corr}(X, Y)$.
    \item \textbf{Nur bei gemeinsamer Normalverteilung äquivalent zu Unabhängigkeit:} $\rho = 0 \Leftrightarrow$ Unabhängigkeit nur für $(X,Y) \sim N(\mu_X, \mu_Y, \sigma_X^2, \sigma_Y^2, \rho)$. In anderen Fällen: Unkorreliertheit $\not\Rightarrow$ Unabhängigkeit.
\end{itemize}

\textbf{Korrelationsmatrix:} Für mehrdimensionale Vektoren $(X_1, \ldots, X_p)$ die $p \times p$-Matrix:
\formel{\mathbf{R} = \begin{pmatrix}
        1         & \rho_{12} & \cdots & \rho_{1p} \\
        \rho_{21} & 1         & \cdots & \rho_{2p} \\
        \vdots    & \vdots    & \ddots & \vdots    \\
        \rho_{p1} & \rho_{p2} & \cdots & 1
    \end{pmatrix}}
mit $\rho_{ij} = \rho_{ji}$ und $\rho_{ii} = 1$.

\textbf{Heatmap-Visualisierung:} Korrelationsmatrizen werden oft als Heatmaps gezeigt, um Muster schnell zu erkennen.

\definition{\textbf{Kovarianzmatrix und Varianz-Kovarianz-Matrix:} Für einen Vektor $\mathbf{X} = (X_1, \ldots, X_p)^T$ mit Mittelwert $\boldsymbol{\mu} = (\mu_1, \ldots, \mu_p)^T$:
\formel{\boldsymbol{\Sigma} = \text{Cov}(\mathbf{X}) = \mathbb{E}[(\mathbf{X} - \boldsymbol{\mu})(\mathbf{X} - \boldsymbol{\mu})^T] = \begin{pmatrix}
            \sigma_1^2  & \sigma_{12} & \cdots & \sigma_{1p} \\
            \sigma_{21} & \sigma_2^2  & \cdots & \sigma_{2p} \\
            \vdots      & \vdots      & \ddots & \vdots      \\
            \sigma_{p1} & \sigma_{p2} & \cdots & \sigma_p^2
        \end{pmatrix}}
Die Diagonalelemente sind Varianzen, die Off-Diagonalelemente sind Kovarianzen. Diese Matrix ist symmetrisch und positiv semi-definit.}

\textbf{Korrelation vs. Kausalität — das wichtigste Konzept:}

\textbf{Fundamentale Warnung:} ``Korrelation impliziert NICHT Kausalität!'' Eine hohe Korrelation $\rho \approx 1$ zwischen zwei Variablen bedeutet nicht, dass eine die andere verursacht. Es gibt mehrere Szenarien:
\begin{itemize}[leftmargin=1.5em]
    \item \textbf{Kausalität:} $X$ beeinflusst $Y$ (z.\,B.\ Rauchen $\to$ Lungenkrebs).
    \item \textbf{Umgekehrte Kausalität:} $Y$ beeinflusst $X$.
    \item \textbf{Gemeinsame Ursache:} Beide werden durch eine dritte Variable $Z$ beeinflusst (Confounding). Z.\,B.\ Eiscreme-Konsum und Ertrinken — beide steigen mit der Temperatur.
    \item \textbf{Zufälliger Zusammenhang:} Spurious Correlation, besonders bei vielen Variablen (wenn $p$ groß ist, findet man zufällig hohe Korrelationen).
\end{itemize}

\textbf{Beispiel (Simpson's Paradoxon):} Gesamt-Korrelation kann unterschiedliches Vorzeichen haben als die Korrelation in Subgruppen!

\textbf{Fazit:} Für Kausalitätsaussagen braucht man experimentelle Designs oder Kausalitäts-Modellierung (Causal Inference), nicht nur Korrelationen.
\newpage

\subsection{Regressionen}

\definition{\textbf{Ziel:} Modelliere Zusammenhang zwischen Prädiktoren $x_1, \ldots, x_p$ und Response $y$. Anwendungen: Vorhersage • Inferenz über Effekte • Modellierung funktionaler Zusammenhänge $y=f(\mathbf{x})+\varepsilon$.}

\subsubsection{Einfache Lineare Regression (SLR)}

\definition{\textbf{Modell:} $y_i = \beta_0 + \beta_1 x_i + \varepsilon_i$ mit $\beta_0$ (Intercept), $\beta_1$ (Slope), $\varepsilon_i \sim N(0, \sigma^2)$. \textbf{Annahmen:} Linearität • Unabhängigkeit • Homoskedastizität (konstante Varianz) • Normalverteilung Fehler.}

\textbf{OLS-Schätzer:} $\hat{\beta}_1 = \frac{s_{xy}}{s_x^2} = r_{xy}\frac{s_y}{s_x}$; $\hat{\beta}_0 = \bar{y} - \hat{\beta}_1 \bar{x}$; Vorhersage $\hat{y} = \hat{\beta}_0 + \hat{\beta}_1 x$.

\textbf{Varianzzerlegung:} TSS $= \sum(y_i - \bar{y})^2 = $ ESS $+ $ RSS mit ESS $= \sum(\hat{y}_i - \bar{y})^2$ (erklärt).

\textbf{R²:} $R^2 = \frac{\text{ESS}}{\text{TSS}} = 1 - \frac{\text{RSS}}{\text{TSS}} \in [0,1]$ = Anteil der Variabilität erklärt. Für SLR: $R^2 = r_{xy}^2$.

\subsubsection{Multiple Lineare Regression (MLR)}

\definition{\textbf{Modell:} $y_i = \beta_0 + \beta_1 x_{i1} + \cdots + \beta_p x_{ip} + \varepsilon_i$. Matrixform: $\mathbf{y} = \mathbf{X}\boldsymbol{\beta} + \boldsymbol{\varepsilon}$.}

\textbf{OLS:} $\hat{\boldsymbol{\beta}} = (\mathbf{X}^T\mathbf{X})^{-1}\mathbf{X}^T\mathbf{y}$; $\hat{\mathbf{y}} = \mathbf{X}\hat{\boldsymbol{\beta}}$.

\textbf{Gütemaße:} $R^2$ (steigt immer, auch mit unwichtigen Prädiktoren) • $R^2_{\text{adj}} = 1 - \frac{(1-R^2)(n-1)}{n-p-1}$ (bevorzugt) • RSE $= \sqrt{\frac{\text{RSS}}{n-p-1}}$.

\textbf{Inferenz:} t-Tests für einzelne $\beta_j$ • F-Test für Gesamtmodell.

\subsubsection{Vergleich: SLR vs. MLR vs. Multivariat}

\textbf{Multivariat:} Mehrere Response-Variablen $y_1, \ldots, y_k$; Modell: $\mathbf{Y} = \mathbf{X}\mathbf{B} + \mathbf{E}$.
\begin{center}
    \small
    \begin{tabular}{|l|l|l|l|}
        \hline
        \textbf{Typ} & \textbf{Prädiktoren} & \textbf{Response} & \textbf{Koeffiziente}                               \\
        \hline
        SLR          & $1 (x)$               & $1 (y)$           & $\beta_0, \beta_1$                                  \\
        \hline
        MLR          & $p \geq 2$           & $1 (y)$           & $\boldsymbol{\beta} = (\beta_0, \ldots, \beta_p)^T$ \\
        \hline
        Multivariat  & $p \geq 1$           & $k \geq 2$        & Matrix $\mathbf{B}$ $(p+1) \times k$                \\
        \hline
    \end{tabular}
\end{center}

\subsubsection{Logistische Regression}

\definition{\textbf{Zweck:} Regression für binäre Response $y \in \{0,1\}$ (Klassifikation). \textbf{Problem linear:} Vorhersagen außerhalb $[0,1]$, nicht konstante Fehlervarianz.}

\definition{\textbf{Modell:} Modelliert $p_i = P(y_i=1|\mathbf{x}_i)$ via Sigmoid (logistische Funktion):
$$p_i = \frac{1}{1+e^{-(\beta_0 + \boldsymbol{\beta}^T\mathbf{x}_i)}}$$
Log-Odds (Logit, linear): $\log\left(\frac{p_i}{1-p_i}\right) = \beta_0 + \beta_1 x_{i1} + \cdots + \beta_p x_{ip}$.}

\textbf{Koeffizienteninterpretation:} Anstieg $x_j$ um 1 multipliziert Odds mit $e^{\beta_j}$.

\textbf{Schätzung:} Maximum Likelihood (nicht OLS), iterativ (Newton-Raphson).

\textbf{Gütemaße:} Deviance • McFadden $R^2$ (Pseudo) • Confusion Matrix (Accuracy, Sensitivity, Specificity) • ROC-AUC.
\newpage




\subsection{Polynomregression, Splines und additive Modelle}

\definition{\textbf{Polynomregression:} $y_i = \beta_0 + \beta_1 x_i + \beta_2 x_i^2 + \cdots + \beta_d x_i^d + \varepsilon_i$ (Grad $d$). Reformuliert als MLR mit $x_{ij} = x_i^j$; geschätzt via OLS. \textbf{Problem:} Höheres $d$ führt zu Overfitting (hohe Varianz). Wahl $d$ via Cross-Validation, AIC, BIC. Nachteil: Global polynomiale Struktur (wenig flexibel bei lokalen Abweichungen).}

\definition{\textbf{Splines (Cubic Spline):} Stückweise polynomiale Funktionen. Teile $x$-Bereich in $K$ Intervalle (Knots), fit lokal kubisches Polynom in jedem, mit Stetigkeit \& Glattheits-Bedingungen (Sekund-Ableitungen stetig) an Knots. \textbf{Vorteile:} Flexibel lokal • weniger Overfitting als Polynome • Knot-Wahl wichtig (Cross-Validation). \textbf{Natural Cubic Spline:} Zusätzliche Bedingung: Außerhalb äußerer Knots ist Funktion linear (stabilisiert Extrapolation).}

\definition{\textbf{Smoothing Spline:} Hybrid Spline-Regression \& Glättung. Minimiert:
$$\sum_i (y_i - f(x_i))^2 + \lambda \int [f''(x)]^2 dx$$
Erster Term: Anpassung; zweiter Term ($\lambda$ Regularisierungsparameter): Bestraft Rauhheit. Größeres $\lambda$ = glatter. $\lambda$ via Cross-Validation (GCV oder LOOCV).}

\definition{\textbf{Generalized Additive Models (GAM):} Nichtlineare Regression mit separaten Funktionen pro Prädiktor:
$$y_i = \beta_0 + f_1(x_{i1}) + \cdots + f_p(x_{ip}) + \varepsilon_i$$
Jede $f_j$ ist Smoothing-Funktion (Spline o.ä.). \textbf{Vorteile:} Interpretierbar • lokal flexibel. Estimation: Backfitting (iteratives Glätten).}

\textbf{Vergleich:} Polynomregression (einfach, global) • Splines (flexibel, lokal) • GAM (sehr flexibel, mehrere Prädiktoren). Alle erfordern Regularisierung gegen Overfitting.

\textbf{Anzahl der Freiheitsgrade:} Hängt von der Anzahl der Knots $K$ ab. Mit $K$ Knots hat ein kubischer Spline etwa $K + 3$ Freiheitsgrade (näherungsweise).

\textbf{Mathematische Formulierung via Basis-Funktionen:} Splines können als Linearkombination von Basis-Funktionen dargestellt werden:
\formel{f(x) = \sum_{j=1}^M \beta_j B_j(x)}
wobei $B_j(x)$ Basis-Funktionen sind (z.\,B.\ B-Spline-Basis) und $M$ die Anzahl der Basis-Funktionen ist. Das Modell ist wieder linear in den Koeffizienten und kann mit OLS geschätzt werden!

\textbf{Vergleich: Polynomregression vs.\ Splines:}
\begin{table}[h]
    \centering
    \small
    \begin{tabular}{|l|l|l|}
        \hline
        \textbf{Aspekt}    & \textbf{Polynomregression} & \textbf{Splines}                \\
        \hline
        Flexibilität       & Global (hängt von $d$ ab)  & Lokal (variabel in Intervallen) \\
        \hline
        Overfitting-Risiko & Höher bei großem $d$       & Besser kontrollierbar           \\
        \hline
        Glattheit          & Automatisch Grad $d-1$     & Kontrollierbar (typisch C2)     \\
        \hline
        Extrapolation      & Polynomial explodiert      & Oft linear außerhalb            \\
        \hline
        Interpretation     & Einfacher                  & Komplexer                       \\
        \hline
    \end{tabular}
\end{table}

\definition{\textbf{Smoothing Splines:} Eine regularisierte Variante von Splines, die einen Kompromiss zwischen Anpassung und Glattheit findet. \textbf{Ziel-Funktion:} Minimiere:
\formel{\sum_{i=1}^n (y_i - f(x_i))^2 + \lambda \int (f''(x))^2 dx}}
wobei:
\begin{itemize}[leftmargin=1.5em]
    \item Erster Term: Anpassung an Daten (wie OLS).
    \item Zweiter Term: Glattheits-Strafterm (je rauer $f$, desto höher der Strafterm).
    \item $\lambda \geq 0$: Regularisierungsparameter (Glättungskonstante).
\end{itemize}
Bei $\lambda = 0$: Perfekte Interpolation (sehr wiggly). Bei $\lambda \to \infty$: Lineare Regressionsgerade. Die optimale $\lambda$ wird oft via \textbf{Cross-Validation} gewählt.

\subsubsection{Additive Modelle (GAM)}

\definition{\textbf{Generalized Additive Models (GAM):} Ein flexibles Regressions-Framework, das additive Struktur nutzt:
\formel{y_i = \beta_0 + f_1(x_{i1}) + f_2(x_{i2}) + \cdots + f_p(x_{ip}) + \varepsilon_i}
wobei $f_j$ \textbf{glatte Funktionen} sind (nicht unbedingt linear). Der große Vorteil: Man schätzt jede Funktion $f_j$ \textbf{separat}, eine Variable nach der anderen.}

\textbf{Flexibilität vs.\ Interpretierbarkeit:}
\begin{itemize}[leftmargin=1.5em]
    \item \textbf{Lineares Modell:} $f_j(x_j) = \beta_j x_j$ (starr, aber interpretierbar).
    \item \textbf{Additive Modelle:} $f_j$ sind glatte Funktionen (flexibel, aber weniger interpretierbar).
    \item \textbf{Allgemeines nichtlineares Modell:} $y = f(\mathbf{x}) + \varepsilon$ (höchste Flexibilität, schwer zu verstehen).
\end{itemize}

\definition{\textbf{Schätzung in GAM via Backfitting:} Ein iterativer Algorithmus:}
\begin{enumerate}[leftmargin=2em]
    \item Initialisiere $\hat{\beta}_0 = \bar{y}$.
    \item Für jede Variable $j = 1, \ldots, p$:
          \begin{itemize}
              \item Berechne das Residuum, ``abzüglich'' aller anderen Funktionen: $\mathbf{r}_j = \mathbf{y} - \hat{\beta}_0 - \sum_{k \neq j} \hat{f}_k(\mathbf{x}_k)$.
              \item Schätze $\hat{f}_j$ durch Glätten von $\mathbf{r}_j$ gegen $\mathbf{x}_j$ (z.\,B.\ via Smoothing Spline).
          \end{itemize}
    \item Wiederhole, bis Konvergenz.
\end{enumerate}
Dies ist ein \textbf{greedy} Algorithmus, der schnell ist, aber nicht garantiert ein globales Optimum findet.

\textbf{Basis-Funktionen in GAM:} Jede $f_j$ wird als Linearkombination von Basis-Funktionen dargestellt:
\formel{f_j(x_j) = \sum_{k=1}^{K_j} \beta_{jk} B_{jk}(x_j)}
Beliebte Wahlen: B-Splines, natural cubic splines, smooth functions. Mit dieser Darstellung wird das GAM-Optimization wieder zu einem \textbf{regularisierten Linearen Problem}.

\definition{\textbf{Variablen-Auswahl in GAM:}}
\begin{itemize}[leftmargin=1.5em]
    \item \textbf{Partial Dependence Plot:} Plot von $\hat{f}_j(x_j)$ gegen $x_j$, zeigt den isolierten Effekt von $x_j$.
    \item \textbf{Effectivity:} ANOVA-ähnliche Tests, um zu prüfen, ob $f_j$ signifikant ist (gegen $f_j = \text{const}$).
    \item \textbf{Regularisierung:} Penalty auf Rauheit zwingt unwichtige Variablen zu $f_j(x_j) = \text{const}$ (automatische Selection).
\end{itemize}

\textbf{GAM vs.\ Lineare Regression:}
\begin{table}[h]
    \centering
    \small
    \begin{tabular}{|l|l|l|}
        \hline
        \textbf{Aspekt}     & \textbf{Lineare Regression}                    & \textbf{GAM}                                \\
        \hline
        Modell              & $y = \beta_0 + \sum \beta_j x_j + \varepsilon$ & $y = \beta_0 + \sum f_j(x_j) + \varepsilon$ \\
        \hline
        Flexibilität        & Niedrig (linear)                               & Höher (nichtlinear)                         \\
        \hline
        Interpretierbarkeit & Hoch (Koeff. klar)                             & Moderat (Plots nötig)                       \\
        \hline
        Rechenkomplexität   & Niedrig (geschlossen)                          & Höher (iterativ)                            \\
        \hline
        Overfitting         & Geringer                                       & Höher (erfordert Regularisierung)           \\
        \hline
    \end{tabular}
\end{table}

\definition{\textbf{Multivariate GAM:} Erweiterung auf \textbf{mehrere Response-Variablen}:
\formel{y_{ij} = \beta_{0j} + f_{1j}(x_{i1}) + \cdots + f_{pj}(x_{ip}) + \varepsilon_{ij} \quad (j = 1, \ldots, k)}
oder mit geteilten Basis-Funktionen für die \textbf{simultane} Schätzung aller Responses. Dies ist besonders wertvoll, wenn die Responses korreliert sind.}

\textbf{Praktische Anwendungen:}
\begin{itemize}[leftmargin=1.5em]
    \item \textbf{Polynomregression:} Wenn die Beziehung quadratisch oder kubisch vermutet wird; einfach zu interpretieren.
    \item \textbf{Splines:} Für glatte, nichtlineare Beziehungen mit lokal variierender Komplexität.
    \item \textbf{GAM:} Für viele Prädiktoren mit komplexen Abhängigkeiten, wenn Interpretierbarkeit wichtig ist (Kompromiss zwischen linearer Regression und Black-Box-Modellen).
\end{itemize}

\newpage

\subsection{PCA, PCR und MDS}

\definition{\textbf{PCA (Principal Component Analysis):} Unsupervisierte Dimensionsreduktion. Findet neue orthogonale Koordinaten (PCs), bei denen maximale Varianz in ersten wenigen Achsen konzentriert ist.}

\textbf{Kern:} Datenmatrix $\mathbf{X}$ ($n \times p$) → Neue Koordinaten: $\mathbf{X} = \mathbf{T}\mathbf{P}^T + \mathbf{E}$ mit Score-Matrix $\mathbf{T}$ (unkorreliert) und Loading-Matrix $\mathbf{P}$ (Eigenvektoren).

\textbf{Algorithmus:} (1) Standardisierung $\mathbf{X}$ • (2) Kovarianzmatrix $\mathbf{C} = \frac{1}{n-1}\mathbf{X}^T\mathbf{X}$ • (3) Eigenwertzerlegung: $\lambda_1 \geq \lambda_2 \geq \cdots$ • (4) Scores: $\mathbf{T} = \mathbf{X}\mathbf{P}$.

\textbf{Varianzzerlegung:} Kumulativer Varianzanteil der ersten $a$ PCs: $\frac{\sum_{j=1}^a \lambda_j}{\sum_{j=1}^p \lambda_j}$. Scree-Plot zur Wahl $a$ (Ellbogen).

\textbf{Visualisierung:} Score-Plot (PC1 vs PC2 der Beobachtungen) • Loading-Plot (welche Originalvariablen beeinflussen jede PC) • Biplot (kombiniert beide).

\textbf{Anwendungen:} Visualisierung • Datenreduktion • Identifikation von Clustern/Ausreißern • Rausch-Reduktion.

\definition{\textbf{PCR (Principal Component Regression):} Löst Multikollinearität \& $p > n$ Problem.}

Nutzt $a$ PCA-Scores statt Originalvariablen als Prädiktoren:
$$\mathbf{y} = \mathbf{T}\boldsymbol{\vartheta} + \boldsymbol{\varepsilon}$$
Vorteil: Scores sind orthogonal (unkorreliert) → OLS stabil. $a$ via Cross-Validation.

\definition{\textbf{MDS (Multidimensional Scaling):} Dimensionsreduktion aus Distanzmatrix (nicht Rohdaten). Input: $n \times n$ Distanzmatrix $\mathbf{D}$ mit $d_{ij}$ = Abstand $i$ zu $j$. Ziel: Finde niedrigdimensionale Koordinaten $\mathbf{Z}$ mit Distanzen $\approx$ ursprüngliche.}

\textbf{Varianten:} Metric MDS (reproduziere absolute Distanzen) • Non-Metric MDS (reproduziere nur Ordnung, robuster gegen Fehler).

\textbf{Klassisches MDS:} Double Centering $\mathbf{D}$ → Gram-Matrix → Eigenwertzerlegung → Koordinaten.

\textbf{Vergleich:} PCA (Input: Daten) vs MDS (Input: Distanzen). PCA max. Varianz; MDS reproduziert Distanzen.

\begin{itemize}[leftmargin=1.5em]
    \item \textbf{PCA:} Bioinformatik (Gen-Expression), Bildkompression, Feature Engineering für maschinelles Lernen.
    \item \textbf{PCR:} Hochdimensionale Prädiktoren mit Multikollinearität (z.\,B.\ Spektroskopie-Daten, Gene).
    \item \textbf{MDS:} Psychometrie (Subjektive Ähnlichkeitsurteile), Genetische Distanzen, Chemische Ähnlichkeitsmaße.
\end{itemize}

\newpage

\subsection{Lasso und Ridge-Regression}

\definition{\textbf{Regularisierung:} Techniken zur Begrenzung von Koeffizientengrößen um Overfitting zu vermeiden bei Multikollinearität oder hochdimensionalen Daten.}

\definition{\textbf{Ridge Regression:} Minimiere:
$$L_{\text{ridge}} = \|\mathbf{y} - \mathbf{X}\boldsymbol{\beta}\|^2_2 + \lambda \sum_j \beta_j^2 = \text{RSS} + \lambda \|\boldsymbol{\beta}\|_2^2$$
Geschlossene Lösung: $\hat{\boldsymbol{\beta}}_{\text{ridge}} = (\mathbf{X}^T\mathbf{X} + \lambda \mathbf{I})^{-1}\mathbf{X}^T\mathbf{y}$. Effekt: Alle Koeffizienten werden geschrumpft gegen Null (nie exakt Null). Vorteil: Stabil bei Multikollinearität.}

\definition{\textbf{Lasso (L1-Penalty):} Minimiere:
$$L_{\text{lasso}} = \|\mathbf{y} - \mathbf{X}\boldsymbol{\beta}\|^2_2 + \lambda \sum_j |\beta_j| = \text{RSS} + \lambda \|\boldsymbol{\beta}\|_1$$
Keine geschlossene Lösung (numerisch via Coordinate Descent). Kritischer Unterschied: \textbf{Setzt Koeffizienten exakt zu Null} $\to$ Variablen-Auswahl. Geometrisch: Diamond-Shape Constraint (Ecken auf Achsen) vs. Ridge-Circle.}

\begin{center}
    \small
    \begin{tabular}{|l|l|l|}
        \hline
        \textbf{Eigenschaft} & \textbf{Ridge (L2)} & \textbf{Lasso (L1)} \\
        \hline
        Penalty              & $\sum \beta_j^2$    & $\sum |\beta_j|$    \\
        \hline
        Lösung               & Geschlossen         & Numerisch           \\
        \hline
        Koeff. $= 0$?        & Selten              & Oft                 \\
        \hline
        Variablen-Auswahl    & Nein                & Ja (Sparsity)       \\
        \hline
        Multikollinearität   & Gut                 & Wählt eine aus      \\
        \hline
        Best für             & Alle Var. relevant  & Sparse, $p \gg n$   \\
        \hline
    \end{tabular}
\end{center}

\definition{\textbf{Elastic Net:} Kombiniert Ridge + Lasso mit $\alpha \in [0,1]$ Mixing-Parameter:
$$L_{\text{EN}} = \text{RSS} + \lambda[\alpha \|\boldsymbol{\beta}\|_1 + (1-\alpha) \|\boldsymbol{\beta}\|_2^2]$$
$\alpha=0$ reine Ridge; $\alpha=1$ reines Lasso; Kompromiss für Korrelation + Sparsity.}

\textbf{Praktisch:} $\lambda$ \& $\alpha$ via Cross-Validation wählen. Ridge nutzen wenn alle Vars relevant; Lasso wenn Variablen-Auswahl gewünscht; Elastic Net als Kompromiss.

\newpage

\subsection{Entscheidungsbäume (Decision Trees) und Random Forests}

\subsubsection{Decision Trees (Entscheidungsbäume)}

\definition{\textbf{Decision Tree:} Supervisierter Lernalgorithmus für Klassifikation und Regression, der Daten rekursiv in homogenere Teilmengen aufteilt. Die Grundstruktur besteht aus Knoten (Bedingungen $x_j \leq c$), Ästen (ja/nein-Ausgänge) und Blättern (Klassen- oder Vorhersagewerte). Ein neuer Datenpunkt traversiert den Baum von der Wurzel bis zu einem Blatt.}

\definition{\textbf{Baumaufbau -- Recursive Partitioning (Top-Down):}}
\begin{enumerate}[leftmargin=2em]
    \item Starte mit allen Trainingsdaten an der Wurzel.
    \item Wähle Variable $x_j$ und Schwellwert $c$, der die reinsten Kindknoten erzeugt.
    \item Teile in $x_j \leq c$ (links) und $x_j > c$ (rechts); wiederhole rekursiv.
    \item Stopp wenn: alle Punkte einer Klasse angehören, minimale Knotengröße, maximale Tiefe, oder keine weitere Verbesserung möglich.
\end{enumerate}

\definition{\textbf{Unreinheitsmaße:}}

\textbf{Klassifikation:} Der \textbf{Gini-Index} $\text{Gini} = 1 - \sum_{k} p_k^2$ misst die Wahrscheinlichkeit einer Fehlklassifikation (0 = rein). Die \textbf{Entropie} $\text{Entropy} = -\sum_{k} p_k \log_2 p_k$ quantifiziert Unordnung; der \textbf{Information Gain} misst die Entropiereduktion durch eine Aufteilung:
\formel{\text{IG} = \text{Entropy(Parent)} - \sum_{i} \tfrac{n_i}{n}\,\text{Entropy(Child}_i\text{)}}

\textbf{Regression:} Minimierung der Residualsumme der Quadrate $\text{RSS} = \sum_i (y_i - \hat{y})^2$ über beide Kindknoten.

\definition{\textbf{Overfitting und Pruning:} Tiefe Bäume fitten Trainingsdaten perfekt, generalisieren aber schlecht. Gegenmaßnahmen: \textbf{Early Stopping} (Tiefenbegrenzung, Mindestgröße pro Knoten) oder \textbf{Post-Pruning} (voller Baum trainieren, dann Knoten entfernen wenn Validierungsperformance steigt). \textbf{Cost Complexity Pruning} minimiert Trainingsfehler plus Komplexitätsstrafe gemeinsam.}

\begin{table}[h]
    \centering\small
    \begin{tabular}{ll}
        \toprule
        \textbf{Vorteile} & \textbf{Nachteile} \\
        \midrule
        Interpretierbar, visualisierbar & Instabil (kleine Datenänderungen) \\
        Keine Skalierung nötig & Overfitting ohne Regularisierung \\
        Kategoriale + numerische Prädiktoren & Greedy, nicht global optimal \\
        Erfasst Nichtlinearität und Interaktionen & Bias bei unbalancierten Klassen \\
        \bottomrule
    \end{tabular}
\end{table}

\subsubsection{Random Forests}

\definition{\textbf{Random Forest:} Ensemble-Methode, die $B$ Entscheidungsbäume kombiniert. Zwei Schlüsselmechanismen reduzieren Varianz und Korrelation zwischen Bäumen: \textbf{Bootstrap-Sampling} (jeder Baum auf einer Zufallsstichprobe mit Zurücklegen) und \textbf{Random Feature Selection} (bei jedem Split nur $m \approx \sqrt{p}$ zufällig gewählte Prädiktoren). Vorhersage via Mehrheitsvote (Klassifikation) oder Durchschnitt (Regression).}

\definition{\textbf{Out-of-Bag (OOB) Error:} Da Bootstrap-Stichproben ca.\ $37\%$ der Daten auslassen, können diese zur Validierung genutzt werden -- ohne separaten Validierungssatz. Jede Beobachtung wird nur durch Bäume vorhergesagt, in deren Stichprobe sie nicht enthalten war.}

\definition{\textbf{Feature Importance:} Zwei Ansätze: \textbf{Mean Decrease in Impurity (MDI)} -- summiere über alle Bäume die Impurity-Reduktion bei Splits auf Variable $X_j$; größere Reduktion = wichtigere Variable. \textbf{Permutation Importance (MDA)} -- permutiere $X_j$ zufällig und messe den OOB-Fehleranstieg; stärkere Verschlechterung = wichtigere Variable.}

\begin{table}[h]
    \centering\small
    \begin{tabular}{lll}
        \toprule
        \textbf{Aspekt} & \textbf{Decision Tree} & \textbf{Random Forest} \\
        \midrule
        Modell & Einzelner Baum & Ensemble von $B$ Bäumen \\
        Varianz & Hoch & Niedrig \\
        Interpretierbarkeit & Hoch & Niedrig (Black Box) \\
        Performance & Moderat & Höher \\
        Rechenzeit & Niedrig & Höher \\
        Feature Importance & Einfach & MDI oder Permutation \\
        \bottomrule
    \end{tabular}
\end{table}

\definition{\textbf{Weitere Ensemble-Methoden:} \textbf{Bagging} verallgemeinert Bootstrap-Aggregation auf beliebige Lerner. \textbf{Boosting} (AdaBoost, Gradient Boosting) trainiert Bäume sequenziell, wobei jeder Baum die Fehler des vorigen korrigiert -- höhere Accuracy, aber overfitting-anfälliger. \textbf{XGBoost} ist eine regularisierte, hochoptimierte GBM-Implementierung, die in Praxis und Wettbewerben weit verbreitet ist.}

\newpage

\subsection{Linear Discriminant Analysis (LDA) und Quadratic Discriminant Analysis (QDA)}

\subsubsection{Klassifikation via Wahrscheinlichkeitsdichten}

\definition{\textbf{Setting:} Gegeben Datenmatrix $\mathbf{X}$ ($n \times p$) und kategoriale Response $y \in \{1,\ldots,K\}$. Beobachtungen in Klasse $k$ entstammen einer Dichte $f_k(\mathbf{x})$. Die \textbf{ML-Klassifikationsregel} ordnet $\mathbf{x}$ der Klasse mit höchster Dichte zu: $\hat{k} = \arg\max_k f_k(\mathbf{x})$. Mit \textbf{Prior-Wahrscheinlichkeiten} $p_k$ wird nach der Bayes-Regel klassifiziert:
\formel{\hat{k} = \arg\max_{k} f_k(\mathbf{x}) \cdot p_k}
Intuition: Selbst wenn $f_2(\mathbf{x}) > f_1(\mathbf{x})$, kann Klasse~1 bevorzugt werden, falls $p_1 \gg p_2$.}

\subsubsection{Linear Discriminant Analysis (LDA)}

\definition{\textbf{LDA-Annahmen:} Jede Klasse $k$ ist multivariat normalverteilt, $\mathbf{x}|C{=}k \sim N_p(\boldsymbol{\mu}_k, \boldsymbol{\Sigma})$, mit \emph{gemeinsamer} Kovarianzmatrix $\boldsymbol{\Sigma}$ -- dies erzeugt \textbf{lineare} Entscheidungsgrenzen. Im multivariaten Fall lautet die \textbf{Diskriminanzfunktion}:
\formel{\delta_k(\mathbf{x}) = \mathbf{x}^T \boldsymbol{\Sigma}^{-1} \boldsymbol{\mu}_k - \tfrac{1}{2}\boldsymbol{\mu}_k^T \boldsymbol{\Sigma}^{-1} \boldsymbol{\mu}_k + \log p_k}
Klassifikation: $\hat{k} = \arg\max_k \delta_k(\mathbf{x})$. Grenzen zwischen Klassen $k$ und $k'$ (wo $\delta_k = \delta_{k'}$) sind lineare Hyperebenen. \textbf{Parameterschätzung:} Klassenmittelwerte $\hat{\boldsymbol{\mu}}_k = \frac{1}{n_k}\sum_{i:y_i=k}\mathbf{x}_i$, Priors $\hat{p}_k = n_k/n$, gepoolte Kovarianzmatrix $\hat{\boldsymbol{\Sigma}}_{\text{pooled}} = \frac{1}{n-K}\sum_k (n_k-1)\hat{\boldsymbol{\Sigma}}_k$.}

\subsubsection{Quadratic Discriminant Analysis (QDA)}

\definition{\textbf{QDA-Annahmen:} Wie LDA, aber jede Klasse hat eine \emph{eigene} Kovarianzmatrix $\boldsymbol{\Sigma}_k$ -- dies erzeugt \textbf{quadratische} Entscheidungsgrenzen (Quadriken). Die \textbf{Diskriminanzfunktion} lautet:
\formel{\delta_k(\mathbf{x}) = -\tfrac{1}{2}\log|\boldsymbol{\Sigma}_k| - \tfrac{1}{2}(\mathbf{x}-\boldsymbol{\mu}_k)^T\boldsymbol{\Sigma}_k^{-1}(\mathbf{x}-\boldsymbol{\mu}_k) + \log p_k}
Der quadratische Term $\mathbf{x}^T\boldsymbol{\Sigma}_k^{-1}\mathbf{x}$ unterscheidet QDA von LDA. Pro Klasse werden $p(p+1)/2$ Kovarianzparameter geschätzt -- bei kleinem $n$ besteht erhöhtes Overfitting-Risiko.}

\begin{table}[h]
    \centering\small
    \begin{tabular}{lll}
        \toprule
        \textbf{Aspekt} & \textbf{LDA} & \textbf{QDA} \\
        \midrule
        Kovarianzannahme & Gemeinsam $\boldsymbol{\Sigma}$ & Klassenspezifisch $\boldsymbol{\Sigma}_k$ \\
        Entscheidungsgrenze & Linear & Quadratisch \\
        Parameter & Wenige & Viele (pro Klasse) \\
        Bias/Varianz & Höherer Bias & Höhere Varianz \\
        Bevorzugt wenn & Ähnliche Spreads, kleines $n$ & Unterschiedliche Spreads, großes $n$ \\
        \bottomrule
    \end{tabular}
\end{table}

\subsubsection{Varianten und verwandte Methoden}

\definition{\textbf{Fisher's Linear Discriminant} (binär) findet die Projektionsrichtung $\mathbf{w} \propto \boldsymbol{\Sigma}_{\text{pooled}}^{-1}(\bar{\mathbf{x}}_1 - \bar{\mathbf{x}}_2)$, die Between-Class- zu Within-Class-Variabilität maximiert -- konzeptuell auf Dimensionsreduktion ausgerichtet, ohne Normalverteilungsannahme.}

\definition{\textbf{Regularized DA (RDA):} Interpoliert zwischen LDA und QDA via $\hat{\boldsymbol{\Sigma}}_k(\lambda) = \lambda\hat{\boldsymbol{\Sigma}}_{\text{pooled}} + (1-\lambda)\hat{\boldsymbol{\Sigma}}_k$, wobei $\lambda \in [0,1]$ per Cross-Validation gewählt wird.}

\definition{\textbf{Mixture DA (MDA):} Modelliert jede Klasse als Mischung mehrerer Normalverteilungen für multimodale Klassenverteilungen.}

\definition{\textbf{Methodenwahl:} LDA bei ähnlichen Kovarianzstrukturen und kleinen Stichproben; QDA bei deutlich unterschiedlichen Spreads und ausreichend Daten ($n \gg p \cdot K$); RDA bei Unsicherheit zwischen beiden; andere Klassifikatoren (SVM, kNN, Trees) bei stark verletzten Normalverteilungsannahmen oder nichtlinearen Grenzen.}

\newpage

\subsection{k-Nearest Neighbors (kNN) -- Klassifikation und Regression}

\definition{\textbf{kNN-Grundkonzept:} Nicht-parametrischer \emph{Lazy Learner} -- es gibt keine Trainingsphase, die Daten werden direkt gespeichert. Für eine neue Beobachtung $\mathbf{x}_0$ werden die $k$ nächsten Nachbarn per Distanzmetrik gesucht; Vorhersage via Mehrheitsvote (Klassifikation) oder Durchschnitt (Regression).}

\subsubsection{Distanzmaße}

\definition{\textbf{Distanzfunktion} muss vier Metrik-Axiome erfüllen: Nicht-Negativität, Identität ($d=0 \Leftrightarrow \mathbf{x}=\mathbf{y}$), Symmetrie und Dreiecksungleichung.}

Gebräuchliche Metriken: \textbf{Euklidisch} $\|\mathbf{x}_A - \mathbf{x}_B\|_2$ (Standard), \textbf{Manhattan} $\|\mathbf{x}_A - \mathbf{x}_B\|_1$, und allgemein die \textbf{Minkowski-Distanz}
\formel{d_p(\mathbf{x}_A,\mathbf{x}_B) = \Bigl(\sum_j |x_{A,j}-x_{B,j}|^p\Bigr)^{1/p}}
mit $p{=}1$ (Manhattan), $p{=}2$ (Euklidisch), $p{\to}\infty$ (Chebyshev, Maximum).

Die \textbf{Mahalanobis-Distanz} $\sqrt{(\mathbf{x}_A-\mathbf{x}_B)^T\boldsymbol{\Sigma}^{-1}(\mathbf{x}_A-\mathbf{x}_B)}$ berücksichtigt Korrelation und Skalierung. Für binäre Vektoren eignet sich die \textbf{Tanimoto-Distanz} $d = 1 - \frac{\mathbf{x}_A^T\mathbf{x}_B}{\|\mathbf{x}_A\|_1 + \|\mathbf{x}_B\|_1 - \mathbf{x}_A^T\mathbf{x}_B}$.

\textbf{Wichtig:} Variablen müssen normalisiert/standardisiert werden, da unterschiedliche Skalen die Distanz dominieren.

\subsubsection{Wahl von $k$ und Entscheidungsgrenzen}

\definition{\textbf{Einfluss von $k$:} Kleines $k$ (z.\,B.\ $k{=}1$) erzeugt flexible, zackige Grenzen mit hohem Overfitting-Risiko; großes $k$ führt zu glatten Grenzen mit Underfitting-Risiko. Der optimale Wert wird via \textbf{Cross-Validation} bestimmt; bei binärer Klassifikation ist ungerades $k$ sinnvoll um Unentschieden zu vermeiden.}

\subsubsection{Vor- und Nachteile sowie Varianten}

\begin{table}[h]
    \centering\small
    \begin{tabular}{ll}
        \toprule
        \textbf{Vorteile} & \textbf{Nachteile} \\
        \midrule
        Einfach, keine Trainingsphase & $O(n \cdot p)$ pro Vorhersage \\
        Nicht-parametrisch, flexibel & Sensibel gegenüber Skalierung \\
        Klassifikation und Regression & Curse of Dimensionality bei großem $p$ \\
        Keine Verteilungsannahmen & Ausreißer nah am Testpunkt schädlich \\
        \bottomrule
    \end{tabular}
\end{table}

\definition{\textbf{Varianten:} \textbf{Gewichtetes kNN} gewichtet Nachbarn nach Distanz ($w_i = 1/d_i$), sodass nähere Nachbarn mehr Einfluss haben. \textbf{Approximate kNN} verwendet KD-Trees oder Ball-Trees zur Beschleunigung bei großen Datensätzen.}




\subsection{Ähnlichkeitsmaße}
\newpage




\subsection{K-means, K-medoids und Fuzzy-Clustering}

\definition{\textbf{K-means:} Partitioniert $n$ Beobachtungen in $K$ disjunkte Cluster durch Minimierung der Within-Cluster-Sum-of-Squares:
\formel{W = \sum_{k=1}^{K}\sum_{i \in C_k} \|\mathbf{x}_i - \mathbf{c}_k\|_2^2}
wobei $\mathbf{c}_k$ das Zentroid von Cluster $k$ ist. Der Algorithmus alterniert zwischen \textbf{Zuordnung} (jede Beobachtung zum nächsten Zentroid) und \textbf{Update} (neues Zentroid als Clustermittelwert) bis Konvergenz. Die Anzahl $K$ muss vorab festgelegt werden.}

\definition{\textbf{Wahl von $K$:} Die \textbf{Elbow-Methode} plottet WCSS gegen $K$ und sucht den ``Ellbogen'' der monoton fallenden Kurve. Die \textbf{Hopkins-Statistik}
\formel{H = \frac{\sum d_U(i)}{\sum d_U(i) + \sum d_W(i)}}
testet ob überhaupt eine Clusterstruktur vorliegt: $H \approx 0{,}5$ bedeutet zufällige Verteilung (kein Clustering sinnvoll), $H \to 1$ starke Clusterstruktur.}

\definition{\textbf{K-medoids:} Wie K-means, aber Clusterzentren müssen \emph{echte Beobachtungen} sein. Das Medoid minimiert die Summe der Abstände zu allen anderen Clusterpunkten. Vorteile: robuster gegenüber Ausreißern und mit beliebigen Distanzmaßen einsetzbar.}

\definition{\textbf{Fuzzy C-Means (FCM):} Statt harter Zuordnung erhält jede Beobachtung $i$ einen Zugehörigkeitsgrad $u_{ij} \in [0,1]$ zu jedem Cluster $j$ (mit $\sum_j u_{ij} = 1$). Minimiert wird:
\formel{J_m = \sum_{i=1}^{n}\sum_{j=1}^{K} u_{ij}^m \|\mathbf{x}_i - \mathbf{c}_j\|^2}
mit Fuzzifikationsparameter $m > 1$ (typisch $m{=}2$). Memberships und Zentroide werden abwechselnd bis Konvergenz aktualisiert.}

\begin{table}[h]
    \centering\small
    \begin{tabular}{llll}
        \toprule
        \textbf{Aspekt} & \textbf{K-means} & \textbf{K-medoids} & \textbf{Fuzzy C-Means} \\
        \midrule
        Zuordnung & Hart & Hart & Weich (graduell) \\
        Zentren & Durchschnitt & Echte Beobachtung & Gewichteter Durchschnitt \\
        Ausreißer & Empfindlich & Robust & Mittel \\
        Distanzmaß & Euklidisch & Beliebig & Euklidisch \\
        Zusatzparameter & -- & -- & Fuzzifikator $m$ \\
        \bottomrule
    \end{tabular}
\end{table}

\subsection{Model-based Clustering}

\definition{\textbf{Grundkonzept:} Die Daten werden als Mischung von $K$ Wahrscheinlichkeitsverteilungen (üblicherweise multivariate Normalverteilungen) modelliert:
\formel{f(\mathbf{x}) = \sum_{k=1}^{K} \pi_k\, f_k(\mathbf{x} \mid \boldsymbol{\theta}_k)}
mit Mischungsgewichten $\pi_k > 0$, $\sum_k \pi_k = 1$, und komponentenspezifischen Parametern $\boldsymbol{\theta}_k = (\boldsymbol{\mu}_k, \boldsymbol{\Sigma}_k)$.}

\definition{\textbf{EM-Algorithmus:} Iterative Maximum-Likelihood-Schätzung der Mischmodellparameter:}
\begin{enumerate}[leftmargin=2em]
    \item \textbf{E-Schritt:} Berechne posteriore Zugehörigkeitswahrscheinlichkeit (ähnlich Fuzzy-Memberships):
    \formel{\tau_{ik} = \frac{\pi_k f_k(\mathbf{x}_i \mid \boldsymbol{\theta}_k)}{\sum_{j} \pi_j f_j(\mathbf{x}_i \mid \boldsymbol{\theta}_j)}}
    \item \textbf{M-Schritt:} Aktualisiere $\pi_k$, $\boldsymbol{\mu}_k$ und $\boldsymbol{\Sigma}_k$ als gewichtete Mittel mit Gewichten $\tau_{ik}$.
    \item Wiederholen bis Log-Likelihood konvergiert.
\end{enumerate}
\textbf{Vorteile:} Probabilistisches Framework, Modellwahl via AIC/BIC. \textbf{Nachteile:} Rechenaufwendiger als K-means, Normalverteilungsannahme nicht immer realistisch.

\newpage

\subsection{Hierarchisches Clustering}

\definition{\textbf{Grundkonzept:} Erzeugt eine Hierarchie von Partitionen statt einer einzigen festen Lösung. \textbf{Agglomerative} Verfahren starten mit $n$ Einzelclustern und fusionieren iterativ die ähnlichsten -- dies ist der gebräuchlichste Ansatz. Divisive Verfahren (umgekehrt) sind rechenintensiver und seltener.}

\definition{\textbf{Linkage-Methoden} definieren den Abstand zwischen zwei Clustern $C_i$ und $C_j$:}
\begin{itemize}[leftmargin=1.5em]
    \item \textbf{Complete:} $\max_{a \in C_i, b \in C_j} d(a,b)$ -- globulare Cluster, kein Chaining.
    \item \textbf{Single:} $\min_{a \in C_i, b \in C_j} d(a,b)$ -- anfällig für Kettenbildung.
    \item \textbf{Average (UPGMA):} Mittlerer Abstand aller Punktpaare -- Kompromiss.
    \item \textbf{Ward:} $\frac{n_i n_j}{n_i+n_j}\|\mathbf{c}_i - \mathbf{c}_j\|^2$ -- minimiert Within-Cluster-Varianz; oft bestes Ergebnis in der Praxis.
\end{itemize}

\definition{\textbf{Dendrogramm:} Baumdiagramm der Fusionshistorie; Blätter sind Einzelbeobachtungen, die Höhe einer Verbindung entspricht dem Fusionsabstand. Horizontale Schnitte liefern Clusterlösungen für beliebiges $K$.

Der \textbf{Agglomerationskoeffizient} $ac = \frac{1}{n}\sum_i (1 - m(i))$, mit $m(i) = d_{i,\text{first}}/d_{i,\text{last}}$, misst die Stärke der Clusterstruktur; Werte nahe 1 deuten auf klare Clusterung hin.

\textbf{Vorteil:} Kein vorab festgelegtes $K$; Dendrogramm ermöglicht visuelle Exploration. \textbf{Nachteil:} $O(n^2)$ Komplexität; Fusionsentscheidungen sind final.}

\newpage

\subsection{Bedingte Wahrscheinlichkeit und Bayes-Theorem}

\definition{\textbf{Bedingte Wahrscheinlichkeit:}
\formel{P(A \mid B) = \frac{P(A \cap B)}{P(B)}, \quad P(B) > 0}
Interpretation: Nach Eintreten von $B$ werden die möglichen Ausgänge auf $B$ eingeschränkt -- $P(A \mid B)$ ist eine Neueinschätzung von $P(A)$ unter dieser Information. \textbf{Produktregel:} $P(A \cap B) = P(A \mid B)\cdot P(B) = P(B \mid A)\cdot P(A)$, verallgemeinert zur Kettenregel für mehrere Ereignisse. \textbf{Satz der totalen Wahrscheinlichkeit:} Für eine Partition $B_1,\ldots,B_K$ des Stichprobenraums gilt $P(A) = \sum_k P(A \mid B_k)\cdot P(B_k)$.}

\definition{\textbf{Satz von Bayes:}
\formel{P(A \mid B) = \frac{P(B \mid A)\cdot P(A)}{P(B)} = \frac{P(B \mid A)\cdot P(A)}{\sum_k P(B \mid A_k)\cdot P(A_k)}}
Die vier Terme: \textbf{Posterior} $P(A|B)$ (aktualisiertes Wissen), \textbf{Likelihood} $P(B|A)$, \textbf{Prior} $P(A)$ (Vorwissen), \textbf{Evidenz} $P(B)$ (Normalisierungskonstante). Der Satz erlaubt ``rückwärts'' zu schließen: von bekanntem $P(B|A)$ auf gesuchtes $P(A|B)$.}

\beispiel{\textbf{Medizinischer Test:} Prävalenz 1\%, Sensitivität 95\%, Spezifität 90\% $\Rightarrow P(+|\text{krank})=0{,}95$, $P(+|\text{gesund})=0{,}10$.
\begin{align*}
P(+) &= 0{,}95 \cdot 0{,}01 + 0{,}10 \cdot 0{,}99 = 0{,}1085 \\
P(\text{krank}\mid+) &= \frac{0{,}95 \cdot 0{,}01}{0{,}1085} \approx 8{,}8\%
\end{align*}
Trotz gutem Test: positives Ergebnis bedeutet nur $\approx 9\%$ Krankheitswahrscheinlichkeit -- die \textbf{Base-Rate Fallacy}.}

\subsection{Probabilistische Graphische Modelle}

\definition{\textbf{Graphische Modelle:} Repräsentation von Wahrscheinlichkeitsverteilungen mittels Graphen, wobei Knoten Zufallsvariablen und Kanten Abhängigkeiten repräsentieren. \textbf{Bayes-Netz:} Ein gerichteter, azyklischer Graph (DAG), bei dem eine Kante $A \to B$ bedeutet, dass $A$ direkter Elternknoten von $B$ ist. Die multivariate Verteilung faktorisiert über die Graphstruktur:
\formel{P(X_1, \ldots, X_n) = \prod_{i=1}^{n} P(X_i \mid \text{Eltern}(X_i))}
Dies reduziert die Parametrisierung erheblich durch bedingte Unabhängigkeiten.}

\definition{\textbf{Bedingte Unabhängigkeit und d-Separation:} Drei kanonische Strukturen bestimmen, wann Informationsfluss blockiert wird:}
\begin{itemize}[leftmargin=1.5em]
    \item \textbf{Seriell} $X \to Y \to Z$: Blockiert wenn $Y$ beobachtet ($X \perp Z \mid Y$).
    \item \textbf{Divergent} $X \leftarrow Y \to Z$: Blockiert wenn $Y$ beobachtet ($X \perp Z \mid Y$).
    \item \textbf{Konvergent} $X \to Y \leftarrow Z$: Blockiert wenn $Y$ \emph{nicht} beobachtet; Beobachten von $Y$ (oder Nachfahren) aktiviert Abhängigkeit -- ``explaining away''.
\end{itemize}

\definition{\textbf{Naive-Bayes-Klassifikator:} Nimmt bedingte Unabhängigkeit aller Features gegeben die Klasse an: $P(X_1, \ldots, X_p \mid C) = \prod_j P(X_j \mid C)$. Klassifikation via MAP-Schätzung, typischerweise im Log-Raum:
\formel{\hat{C} = \arg\max_c \left[\log P(C{=}c) + \sum_{j=1}^{p} \log P(X_j \mid C{=}c)\right]}}

\textbf{Laplace-Smoothing} verhindert Null-Wahrscheinlichkeiten bei ungesehenen Feature-Klasse-Kombinationen:
\formel{P(X_j{=}x \mid C{=}c) = \frac{\text{count}(X_j{=}x,\, C{=}c) + \alpha}{\text{count}(C{=}c) + \alpha \cdot m}}
mit Smoothing-Parameter $\alpha$ (typisch $\alpha=1$) und $m$ möglichen Werten von $X_j$. Vorteil: schnell, benötigt wenig Daten, liefert Klassifikationssicherheiten. Nachteil: Unabhängigkeitsannahme oft unrealistisch, schlechte Kalibrierung bei korrelierten Features.

\newpage

\subsection{Bayes-Updates: Prior, Likelihood, Evidenz, Posterior}

\definition{\textbf{Bayesianische Inferenz:} Unbekannte Parameter werden als Zufallsvariablen behandelt; der Satz von Bayes dient als Update-Regel:
\formel{p(\theta \mid \mathcal{D}) = \frac{p(\mathcal{D} \mid \theta)\cdot p(\theta)}{p(\mathcal{D})} \quad \Leftrightarrow \quad \text{Posterior} = \frac{\text{Likelihood} \times \text{Prior}}{\text{Evidenz}}}}

\begin{itemize}[leftmargin=1.5em]
    \item \textbf{Prior} $p(\theta)$: Vorwissen über $\theta$ vor Beobachtung der Daten (kann uninformativ sein).
    \item \textbf{Likelihood} $p(\mathcal{D}\mid\theta) = \prod_i p(\mathbf{x}_i \mid \theta)$: Wie gut erklärt $\theta$ die Daten?
    \item \textbf{Evidenz} $p(\mathcal{D}) = \int p(\mathcal{D}\mid\theta)\,p(\theta)\,d\theta$: Normalisierungskonstante, wichtig für Modellvergleich.
    \item \textbf{Posterior} $p(\theta\mid\mathcal{D})$: Aktualisiertes Wissen, kombiniert Prior und Likelihood.
\end{itemize}

\definition{\textbf{Sequential Update:} Der alte Posterior wird zum neuen Prior -- Reihenfolge der Datenpunkte ist unerheblich, Batch- und Online-Update liefern identische Ergebnisse.}

\definition{\textbf{Posterior-Punktschätzungen:} \textbf{MAP} $\hat{\theta} = \arg\max_\theta\, p(\theta\mid\mathcal{D})$ (Modus), \textbf{Posterior Mean} $\mathbb{E}[\theta\mid\mathcal{D}]$, \textbf{Posterior Median}. \textbf{Credible Intervals} können direkt probabilistisch interpretiert werden: mit $(1{-}\alpha)$-Kredibilität liegt $\theta$ im Intervall -- im Gegensatz zum frequentistischen Konfidenzintervall.}

\beispiel{\textbf{Münzwurf:} Prior $p \sim \text{Beta}(1,1)$, 5 Würfe mit 3 Erfolgen. Posterior: $\text{Beta}(4,3)$, Erwartungswert $4/7 \approx 0{,}57$. MLE wäre $3/5 = 0{,}6$; der Bayes-Schätzer ist durch den Prior leicht dazu gezogen.}

\subsection{Konjugierte Priors}

\definition{\textbf{Konjugierter Prior:} Ein Prior $p(\theta)$ heißt konjugiert zur Likelihood, wenn der Posterior aus derselben Verteilungsfamilie stammt -- dies erlaubt analytische Update-Regeln ohne numerische Integration.}

\begin{itemize}[leftmargin=1.5em]
    \item \textbf{Beta--Binomial:} Prior $\text{Beta}(\alpha,\beta)$, nach $k$ Erfolgen in $n$ Versuchen: Posterior $\text{Beta}(\alpha+k,\, \beta+n-k)$.
    \item \textbf{Normal--Normal} (bekannte Varianz $\sigma^2$): Prior $N(\mu_0,\tau_0^2)$, Posterior $N(\mu_n,\tau_n^2)$ mit $\tau_n^{-2} = \tau_0^{-2} + n\sigma^{-2}$ und $\mu_n = \tau_n^2(\tau_0^{-2}\mu_0 + n\sigma^{-2}\bar{x})$. Der Posterior ist ein präzisionsgewichteter Durchschnitt von Prior-Mittelwert und Stichprobenmittelwert.
    \item \textbf{Dirichlet--Multinomial:} Prior $\text{Dir}(\alpha_1,\ldots,\alpha_k)$, Posterior $\text{Dir}(\alpha_1+n_1,\ldots,\alpha_k+n_k)$. Anwendung: Naive Bayes, Textklassifikation.
\end{itemize}

\begin{table}[h]
    \centering\small
    \begin{tabular}{lll}
        \toprule
        \textbf{Likelihood} & \textbf{Konjugierter Prior} & \textbf{Posterior} \\
        \midrule
        Bernoulli/Binomial & Beta & Beta \\
        Multinomial & Dirichlet & Dirichlet \\
        Normal ($\mu$, $\sigma^2$ bekannt) & Normal & Normal \\
        Normal ($\sigma^2$, $\mu$ bekannt) & Inverse-Gamma & Inverse-Gamma \\
        Poisson/Exponential ($\lambda$) & Gamma & Gamma \\
        \bottomrule
    \end{tabular}
\end{table}

\definition{\textbf{Normal-Inverse-Gamma (NIW)} -- wenn sowohl $\mu$ als auch $\sigma^2$ unbekannt: Prior $p(\mu,\sigma^2) = N(\mu \mid \mu_0, \sigma^2/\kappa_0)\cdot\text{Inv-Gamma}(\sigma^2 \mid \alpha_0,\beta_0)$. Update-Regeln: $\kappa_n = \kappa_0+n$, $\mu_n = (\kappa_0\mu_0 + n\bar{x})/\kappa_n$, $\alpha_n = \alpha_0+n/2$, $\beta_n = \beta_0 + \frac{1}{2}\sum(x_i-\bar{x})^2 + \frac{n\kappa_0}{2\kappa_n}(\bar{x}-\mu_0)^2$. In der Praxis sehr häufig, wenn keine Annahme über die Varianz getroffen werden soll.}

%-----------------------------------------------------------------------
\newpage






\subsection{Survivorship Bias und kognitive Verzerrungen}

\definition{\textbf{Survivorship Bias:} Analyse basiert nur auf den ``Überlebenden'' eines Prozesses; Misserfolge fehlen systematisch, was zu verzerrten Schlussfolgerungen führt: [cite: 1, 2]}
\formel{\mathbb{E}[\text{Statistik} \mid D_{\text{obs}}] \neq \mathbb{E}[\text{Statistik} \mid D_{\text{obs}} \cup D_{\text{missing}}]}

\beispiel{\textbf{Flugzeug-Panzerfrage (Wald, 2. WK):} Zurückgekehrte Flugzeuge zeigten Treffer an Tragflächen und Rumpf -- naïve Schlussfolgerung: diese Bereiche panzern. Richtige Schlussfolgerung: Flugzeuge, die an Cockpit oder Motoren getroffen wurden, kehrten \emph{nicht} zurück und fehlen in der Analyse. Gerade die Bereiche \emph{ohne} gehäufte Treffer sind die kritischen. [cite: 2, 3, 4]}

\beispiel{\textbf{Helm-Paradoxon (1. WK):} Nach Einführung von Stahlhelmen stiegen registrierte Kopfverletzungen. Erklärung: Vorher starben Soldaten mit Kopftreffern sofort und wurden nie im Feldspital erfasst; Helme erhöhten die Überlebensrate, sodass mehr Verletzte ins Spital kamen. Die scheinbare Zunahme misst Überlebende, nicht Gesamtverletzungen. [cite: 4, 5, 6, 7]}

\definition{\textbf{Availability Bias:} Leicht verfügbare oder medial präsente Ereignisse werden in ihrer Häufigkeit überschätzt (z.\,B.\ Flugzeugunglücke vs.\ Autounfälle). [cite: 7]}

\definition{\textbf{Confirmation Bias:} Nur Daten werden erfasst oder bevorzugt, die bestehende Überzeugungen bestätigen -- in ML-Systemen entsteht dadurch eine gefährliche Rückkopplungsschleife. [cite: 8]}

\beispiel{\textbf{Algorithmischer Confirmation Bias:} Polizei-Einsatzmodell berechnet Verbrechensdichte aus historischen Daten, weist mehr Polizei in betroffene Viertel zu, erfasst dadurch mehr Kleinkriminalität, erhöht die berechnete Dichte weiter -- self-reinforcing loop. Das System verstärkt die historische Verzerrung statt sie zu korrigieren. [cite: 8, 9]}

\definition{\textbf{Gegenmaßnahmen:} Prüfen welche Daten \emph{fehlen} und warum (Selection Mechanism); Missing Data Analysis; Sensitivitätsanalyse; Blind Randomization; bei gesellschaftskritischen ML-Systemen aktives Fairness-Monitoring und Feedback-Loop-Kontrolle. [cite: 9, 10]}

%-----------------------------------------------------------------------
\newpage

\subsection{Wahrheitsmatrix und abgeleitete Kenngrößen}

\definition{\textbf{Confusion Matrix:} Tabellarische Darstellung der Klassifikationsleistung für binäre Klassifikation: [cite: 10]}

\begin{center}
\small
\begin{tabular}{|l|c|c|}
    \hline
    & \textbf{Actual Negative} & \textbf{Actual Positive} \\
    \hline
    \textbf{Predicted Negative} & TN & FN \\
    \hline
    \textbf{Predicted Positive} & FP & TP \\
    \hline
\end{tabular}
\end{center}

FP = Fehlalarm; FN = übersehen. [cite: 10, 11]

\definition{\textbf{Abgeleitete Metriken:}}
\begin{itemize}[leftmargin=1.5em]
    \item \textbf{Accuracy:} $\frac{\text{TP}+\text{TN}}{n}$ -- Gesamtanteil korrekter Vorhersagen; irreführend bei unbalancierten Klassen. [cite: 11]
    \item \textbf{Sensitivity / Recall / TPR:} $\frac{\text{TP}}{\text{TP}+\text{FN}}$ -- Anteil der echten Positiven, die erkannt wurden. [cite: 12]
    \item \textbf{Specificity / TNR:} $\frac{\text{TN}}{\text{TN}+\text{FP}}$ -- Anteil der echten Negativen, die korrekt erkannt wurden. [cite: 13]
    \item \textbf{Precision / PPV:} $\frac{\text{TP}}{\text{TP}+\text{FP}}$ -- Von allen vorhergesagten Positiven: wie viele sind korrekt? [cite: 14]
    \item \textbf{F1-Score:} $\frac{2\cdot\text{Precision}\cdot\text{Recall}}{\text{Precision}+\text{Recall}}$ -- Harmonisches Mittel; bestraft extrem niedrige Werte stärker als arithmetisches Mittel. [cite: 14, 15]
    \item \textbf{Balanced Accuracy:} $\frac{\text{TPR}+\text{TNR}}{2}$ -- Besser als Accuracy bei unbalancierten Daten. [cite: 15, 16]
\end{itemize}

\textbf{Zusammenhang:} Sensitivity und Specificity sind negativ korreliert: höherer Threshold $\Rightarrow$ mehr Spezifität, weniger Sensitivität. Welche bevorzugt wird, hängt vom Kontext ab (Krebsdiagnose: hohe Sensitivity; Spam-Filter: hohe Specificity). [cite: 16, 17]

\definition{\textbf{ROC-Kurve und AUC:} Plottet TPR (Y) gegen FPR $= 1-\text{Specificity}$ (X) über alle Threshold-Werte. Jeder Punkt entspricht einem Threshold; die Diagonale repräsentiert einen Zufallsklassifikator. Je weiter die Kurve nach oben-links gewölbt, desto besser. \textbf{AUC} (Fläche unter der ROC-Kurve) ist ein threshold-unabhängiges Gütemaß: $0{,}5$ = Zufall, $1{,}0$ = perfekt. Wahrscheinlichkeitsinterpretation: AUC = $P(\hat{p}_+ > \hat{p}_-)$ für ein zufälliges positives vs.\ negatives Beispiel. [cite: 17, 18, 19, 20]}

\begin{table}[h]
    \centering\small
    \begin{tabular}{llll}
        \toprule
        \textbf{Metrik} & \textbf{Formel} & \textbf{Fokus} & \textbf{Geeignet für} \\
        \midrule
        Accuracy & $(\text{TP}+\text{TN})/n$ & Gesamtkorrektheit & Balancierte Daten \\
        Sensitivity/Recall & $\text{TP}/(\text{TP}+\text{FN})$ & Positive finden & Medizin, Sicherheit \\
        Specificity & $\text{TN}/(\text{TN}+\text{FP})$ & Negative erkennen & Spam-Filter [cite: 21] \\
        Precision & $\text{TP}/(\text{TP}+\text{FP})$ & Zuverlässigkeit & Kritische Klassifikation \\
        F1-Score & $2PR/(P+R)$ & Balance P \& R & Unbalancierte Daten \\
        AUC & Fläche unter ROC & Ranking-Qualität & Allgemeiner Vergleich \\
        \bottomrule
    \end{tabular}
\end{table}

%-----------------------------------------------------------------------
\newpage

\subsection{Support Vector Machines}

\definition{\textbf{SVM -- Grundidee:} Finde die Hyperebene $H\colon \mathbf{w}^T\mathbf{x} + b = 0$, die zwei Klassen mit maximalem Margin trennt (\emph{Maximum Margin Principle}). Klassifikation: $\hat{y} = \text{sign}(\mathbf{w}^T\mathbf{x}_0 + b)$. Die Lösung hängt nur von den \textbf{Support Vectors} ab -- Trainingspunkten auf oder nahe dem Margin; weiter entfernte Punkte haben keinen Einfluss. [cite: 21, 22, 23]}

\definition{\textbf{Hard-Margin SVM} (linear separierbare Daten):}
\formel{\min_{\mathbf{w},b}\ \tfrac{1}{2}\|\mathbf{w}\|^2 \quad \text{s.t.}\quad y_i(\mathbf{w}^T\mathbf{x}_i + b) \geq 1\ \forall i}
Der Gewichtsvektor ist eine Linearkombination der Trainingspunkte: $\mathbf{w} = \sum_i \alpha_i y_i \mathbf{x}_i$. KKT-Bedingung: $\alpha_i > 0$ genau für Support Vectors. [cite: 23, 24]

\definition{\textbf{Soft-Margin SVM} (Support Vector Classifier, nicht-separierbare Daten): Slack-Variablen $\xi_i \geq 0$ erlauben Constraint-Verletzungen:}
\formel{\min_{\mathbf{w},b,\boldsymbol{\xi}}\ \tfrac{1}{2}\|\mathbf{w}\|^2 + C\sum_i \xi_i \quad \text{s.t.}\quad y_i(\mathbf{w}^T\mathbf{x}_i+b) \geq 1-\xi_i}
$C$ steuert den Bias-Varianz-Tradeoff: großes $C$ $\Rightarrow$ enges Margin, Overfitting-Risiko; kleines $C$ $\Rightarrow$ breites Margin, mehr Fehler toleriert. Optimales $C$ via Cross-Validation. [cite: 24, 25]

\definition{\textbf{Kernel Trick:} Statt expliziter Transformation $\phi(\mathbf{x})$ in höherdimensionalen Raum wird die Kernel-Funktion $K(\mathbf{x}_i, \mathbf{x}_j) = \phi(\mathbf{x}_i)^T\phi(\mathbf{x}_j)$ direkt verwendet -- erzeugt nichtlineare Entscheidungsgrenzen ohne explizite Transformation. Gebräuchliche Kernels: \textbf{Linear} $\mathbf{x}_i^T\mathbf{x}_j$; \textbf{Polynomial} $(1+\mathbf{x}_i^T\mathbf{x}_j)^d$; \textbf{RBF} $\exp(-\gamma\|\mathbf{x}_i-\mathbf{x}_j\|^2)$ (flexibelster, beliebtester Kernel; $\gamma$ steuert Lokalität); \textbf{Sigmoid} $\tanh(\kappa\mathbf{x}_i^T\mathbf{x}_j+\theta)$. [cite: 25, 26, 27]}

\definition{\textbf{Multi-Klassen-SVM:} \textbf{One-vs-One (OvO):} $\binom{K}{2}$ binäre SVMs, Mehrheitsvote. \textbf{One-vs-All (OvA):} $K$ binäre SVMs, höchster Konfidenz-Score gewinnt. [cite: 27]}

\begin{table}[h]
    \centering\small
    \begin{tabular}{llll}
        \toprule
        \textbf{Aspekt} & \textbf{SVM} & \textbf{Log.\ Regression} & \textbf{Random Forest} \\
        \midrule
        Hochdimensional & Sehr gut & Gut & Mäßig \\
        Kernel möglich & Ja & Nein & Nein \\
        Rechenaufwand & Hoch ($O(n^3)$) [cite: 28] & Niedrig & Mittel \\
        Interpretierbarkeit & Niedrig & Hoch & Mittel \\
        Multi-class nativ & Nein & Ja & Ja \\
        \bottomrule
    \end{tabular}
\end{table}

%-----------------------------------------------------------------------
\newpage

\subsection{Validierung, Underfitting und Overfitting}

\definition{\textbf{Bias-Varianz-Zerlegung:} Der erwartete Testfehler zerfällt in drei Komponenten:}
\formel{\text{MSE} = \text{Bias}^2 + \text{Varianz} + \text{Irreduzibler Fehler}}
\textbf{Komponenten:} Bias = systematischer Fehler (Modell zu simpel); Varianz = Empfindlichkeit gegenüber Trainingsdaten (Modell zu komplex); irreduzibles Rauschen = nicht eliminierbar. [cite: 28, 29]

\begin{table}[h]
    \centering\small
    \begin{tabular}{llll}
        \toprule
        \textbf{Situation} & \textbf{Train-Error} & \textbf{Test-Error} & \textbf{Maßnahme} \\
        \midrule
        Underfitting & Hoch & Hoch & Komplexeres Modell, mehr Features \\
        Optimal & Niedrig & Niedrig & -- \\
        Overfitting & Niedrig & Hoch & Regularisierung, mehr Daten [cite: 30] \\
        \bottomrule
    \end{tabular}
\end{table}

\definition{\textbf{Validierungsstrategien:}}

\textbf{Hold-out:} Aufteilung in Trainings- (60--80\%) und Testset (20--40\%). Einfach, aber hohe Varianz der Schätzung bei kleinen Datensätzen. [cite: 30, 31]

\textbf{$k$-Fold Cross-Validation:} Daten werden in $k$ Folds aufgeteilt; jede Fold dient einmal als Testset, der Rest als Trainingsset. CV-Fehler = $\frac{1}{k}\sum_{i=1}^k \text{Error}_i$. Stabiler als Hold-out, nutzt alle Daten. Spezialfall LOOCV ($k=n$) sehr rechenintensiv. Stratified $k$-Fold erhält Klassenverteilung pro Fold. [cite: 31, 32, 33]

\definition{\textbf{Regularisierung:} Reduziert Overfitting durch Bestrafung großer Koeffizienten. \textbf{L2 (Ridge):} $\lambda\sum w_j^2$ -- schrumpft Koeffizienten, setzt keine auf 0. \textbf{L1 (Lasso):} $\lambda\sum|w_j|$ -- erzeugt Sparsität (Feature Selection). \textbf{Elastic Net:} Kombination beider. Hyperparameter $\lambda$ via Cross-Validation wählen. [cite: 33, 34, 35]}

\textbf{Early Stopping:} Training abbrechen sobald Validierungsfehler zu steigen beginnt; Modell mit minimalem Validierungsfehler speichern. [cite: 35, 36]

%-----------------------------------------------------------------------
\newpage

\subsection{Zeitreihenanalyse}

\definition{\textbf{Zeitreihe:} Sequenz von Beobachtungen $(y_t)_{t=1}^T$ in fester zeitlicher Ordnung. Beobachtungen sind zeitlich abhängig (nicht i.i.d.). Wichtige Eigenschaften: feste Ordnung, typischerweise äquidistant, eine Realisierung eines stochastischen Prozesses $(Y_t)$. [cite: 36, 37]}

\definition{\textbf{Stationarität (schwache Form):} Konstanter Mittelwert $E(Y_t)=\mu$, konstante Varianz $\text{Var}(Y_t)=\sigma^2$, und Autokovarianz $\text{Cov}(Y_t,Y_{t+h})$ hängt nur vom Lag $h$ ab, nicht von $t$. Viele Zeitreihenmodelle (ARMA) setzen Stationarität voraus. [cite: 37, 38]}

\definition{\textbf{Autokorrelationsfunktion (ACF):} Empirische Korrelation der Reihe mit ihren Lags:}
\formel{r_\omega = \frac{\sum_{t=\omega+1}^{T}(y_t-\bar{y})(y_{t-\omega}-\bar{y})}{\sum_{t=1}^{T}(y_t-\bar{y})^2}}
Schnell abfallende ACF: Kurzzeitabhängigkeit (stationär). Langsam abfallende ACF: Nicht-stationär. Saisonale Spitzen: periodische Struktur. Die \textbf{PACF} misst Korrelation bei Lag $\omega$ nach Entfernung des Einflusses dazwischenliegender Lags -- hilfreich zur AR-Ordnungsbestimmung. [cite: 38, 39]

\definition{\textbf{Komponentenmodell:} Eine Zeitreihe wird in Trend $\mu_t$, Saisonalität $s_t$ und Restkomponente $\varepsilon_t$ zerlegt. \textbf{Additiv:} $y_t = \mu_t + s_t + \varepsilon_t$ (konstante Saisonamplitude). \textbf{Multiplikativ:} $y_t = \mu_t \cdot s_t \cdot \varepsilon_t$ (proportionale Amplitude; durch Logarithmieren in additives Modell überführbar). Schätzung via \textbf{Simple Moving Average} mit Fensterbreite $2q+1$: $\hat{\mu}_t = \frac{1}{2q+1}\sum_{i=-q}^{q} y_{t+i}$. [cite: 39, 40, 41, 42]}

\definition{\textbf{AR, MA und ARMA-Prozesse:}}
\begin{itemize}[leftmargin=1.5em]
    \item \textbf{AR($p$):} $Y_t = \sum_{k=1}^p \phi_k Y_{t-k} + \varepsilon_t$ -- aktueller Wert als Regression auf $p$ Lags. PACF bricht bei Lag $p$ ab. [cite: 42, 43]
    \item \textbf{MA($q$):} $Y_t = \varepsilon_t + \sum_{k=1}^q \theta_k \varepsilon_{t-k}$ -- gewichtete Summe vergangener Rauschterme. ACF bricht bei Lag $q$ ab. [cite: 43, 44]
    \item \textbf{ARMA($p,q$):} Kombination beider. Voraussetzung: Stationarität. [cite: 44, 45]
\end{itemize}

\definition{\textbf{ARIMA und SARIMA:} \textbf{ARIMA($p,d,q$)} für nicht-stationäre Reihen: $d$-faches Differenzieren erzeugt Stationarität, dann ARMA. \textbf{SARIMA$(p,d,q)\times(P,D,Q)_s$} fügt saisonale Komponenten mit Periode $s$ hinzu. [cite: 45, 46]}

\textbf{Modellidentifikation:} PACF bricht bei Lag $p$ ab $\Rightarrow$ AR($p$); ACF bricht bei Lag $q$ ab $\Rightarrow$ MA($q$). Modellwahl via AIC/BIC; Residuen müssen White Noise sein (Ljung-Box-Test). [cite: 46, 47, 48]

%-----------------------------------------------------------------------
\newpage

\subsection{Importance Sampling und MCMC}

\definition{\textbf{Monte Carlo Integration:} Schätze $E_p[f(x)] = \int f(x)p(x)\,dx$ durch Ziehen von $n$ Stichproben aus $p$: $\hat{E} = \frac{1}{n}\sum_i f(x_i) \to E_p[f(x)]$ (Gesetz der großen Zahlen). Problem: Direktes Sampling aus $p$ oft intraktabel (z.\,B.\ Bayes-Posterior mit unbekannter Normalisierungskonstante). [cite: 48, 49]}

\definition{\textbf{Importance Sampling:} Ziehe stattdessen aus leicht zu samplender \textbf{Proposal-Verteilung} $q(x)$ und korrigiere mit \textbf{Importance Weights} $w(x) = p(x)/q(x)$:}
\formel{E_p[f(x)] = E_q\!\left[f(x)\frac{p(x)}{q(x)}\right] \approx \sum_i \tilde{w}_i f(x_i), \quad \tilde{w}_i = \frac{w_i}{\sum_j w_j}}
Anforderung: $q(x)>0$ überall wo $p(x)>0$. Schlechte Proposal $\Rightarrow$ extrem variable Gewichte $\Rightarrow$ wenige Stichproben dominieren (\emph{Degeneracy}). Die \textbf{Effective Sample Size} $\text{ESS} = (\sum w_i)^2/\sum w_i^2$ misst Effizienz ($\text{ESS}=n$: optimal; $\text{ESS}\ll n$: schlechte Proposal). In hohen Dimensionen versagt Importance Sampling. [cite: 49, 50, 51, 52]

\definition{\textbf{Markov Chain Monte Carlo (MCMC):} Konstruiere eine Markov Chain mit stationärer Verteilung $p(x)$; nach ausreichend langer Laufzeit (Burn-in verwerfen) sind Stichproben approximativ aus $p$. \textbf{Metropolis-Hastings:} Pro Iteration: (1) schlage $x' \sim q(x'\mid x_t)$ vor; (2) akzeptiere mit $\alpha = \min\!\left(1,\frac{p(x')q(x_t\mid x')}{p(x_t)q(x'\mid x_t)}\right)$; (3) sonst bleibe bei $x_t$. Detailed Balance garantiert stationäre Verteilung $p$. \textbf{Gibbs Sampling:} Für multivariate $\mathbf{X}=(X_1,\ldots,X_p)$ sample iterativ aus den vollständig bedingten Verteilungen $p(x_j\mid x_{-j})$ -- Spezialfall von Metropolis-Hastings mit $\alpha=1$. [cite: 52, 53, 54, 55, 56]}

\definition{\textbf{Konvergenzdiagnostik:} \textbf{Trace Plots} (Chain sollte stationär oszillieren); \textbf{ACF der Chain} (schnell abfallend = gutes Mixing); \textbf{Gelman-Rubin $\hat{R}$} (mehrere Chains, $\hat{R}<1{,}1$ deutet auf Konvergenz); \textbf{ESS} (berücksichtigt Autokorrelation). [cite: 56, 57, 58]}

\begin{table}[h]
    \centering\small
    \begin{tabular}{lll}
        \toprule
        \textbf{Aspekt} & \textbf{Importance Sampling} & \textbf{MCMC} \\
        \midrule
        Unabhängige Samples & Ja & Nein (autokorreliert) \\
        Skalierung (hohe Dim.) & Schlecht & Besser \\
        Normalisierungskonstante & Direkt schätzbar & Nicht direkt \\
        Konvergenzdiagnose & Einfach (ESS) & Aufwendig ($\hat{R}$, ACF) [cite: 59] \\
        Geeignet wenn & Niedrige Dim., gute Proposal & Hohe Dim., komplexer Posterior \\
        \bottomrule
    \end{tabular}
\end{table}